% concatenated from sobolev.tex
\documentclass[12pt]{article}

\usepackage[a4paper,margin=1in]{geometry}
\usepackage{amsmath,amssymb,amsfonts,amsthm,mathtools}
\usepackage{mathrsfs}
\usepackage{xcolor}
\usepackage{microtype}
\usepackage{hyperref}

\hypersetup{
  colorlinks=true,
  linkcolor=blue!55!black,
  citecolor=blue!55!black,
  urlcolor=blue!55!black,
  pdfstartview=FitH
}

\theoremstyle{plain}
\newtheorem{theorem}{Theorem}[section]
\newtheorem{lemma}[theorem]{Lemma}
\newtheorem{proposition}[theorem]{Proposition}
\newtheorem{corollary}[theorem]{Corollary}

\theoremstyle{definition}
\newtheorem{definition}[theorem]{Definition}
\newtheorem{example}[theorem]{Example}
\newtheorem{problem}[theorem]{Problem}

\theoremstyle{remark}
\newtheorem{remark}[theorem]{Remark}

\numberwithin{equation}{section}

\newcommand{\nn}{\mathbb{N}}
\newcommand{\rr}{\mathbb{R}}
\newcommand{\Nzero}{\mathbb{N}_{0}}
\newcommand{\Ls}[1]{\mathcal{L}_{s}^{#1}}

\setcounter{tocdepth}{2}
\emergencystretch=2em

\begin{document}

\title{On the equivalence of Sobolev norms in infinite dimensions}
\author{Zhouzhe Wang\footnote{School of Mathematics,
Sichuan University, Chengdu, 610064, China. E-mail address:
wangzhouzhe@stu.scu.edu.cn.}}
\date{\today}
\maketitle

\begin{abstract}
We prove dimension-free higher-order Sobolev norm estimates on open convex subsets of $\mathbb{R}^n$ with respect to Gaussian measure and use them to obtain norm equivalence on nonempty open convex subsets of $\ell^2$ endowed with a nondegenerate Gaussian measure. To the best of our knowledge, this is the first such equivalence theorem on a proper open subset of an infinite-dimensional Hilbert space, beyond the earlier whole-space results. We also prove the Malliavin--Sobolev norm equivalence for all $p\in[1,\infty)$ and $k\ge2$, including the case $p=1$, $k\ge3$ left open by Addona--Muratori--Rossi in \cite{AddonaMuratoriRossi}.




	
\end{abstract}

\tableofcontents

\section{Introduction}

For an integer $m\geqslant1$, the classical Sobolev norm records a function together with all of its weak derivatives of orders up to $m$.  In many situations the intermediate derivatives are not independent: they can be estimated in terms of the function itself and its derivative of highest order.  The full Sobolev norm is then equivalent to the graph norm of the highest-order differential operator.  This elementary-looking principle underlies interpolation estimates, coercive inequalities, regularity theory, and the identification of operator domains; see, for example, \cite{AdamsFournier}.

In infinite dimensions the corresponding question has to be formulated with respect to a reference measure.  There is no nontrivial locally finite translation-invariant measure on an infinite-dimensional Banach space, and hence the usual Lebesgue-based definition has no direct analogue.  Gaussian measures provide the standard replacement.  The resulting Sobolev theory is closely related to Gaussian analysis, Dirichlet forms, the Ornstein--Uhlenbeck semigroup, stochastic analysis, and Malliavin calculus; see \cite{BogachevGaussian,BogachevSurvey,DP,Nualart2006,ShigekawaBook,Ustunel1995}.

Two distinct questions are often grouped under the heading of Sobolev equivalence in infinite dimensions.  The first concerns the underlying space: one may define Sobolev spaces by completing smooth cylindrical functions, or by requiring weak derivatives through Gaussian integration-by-parts identities, and ask whether the two constructions agree.  This is the infinite-dimensional analogue of the classical ``$H=W$'' problem.  The second question concerns two norms on an already defined Sobolev space.  If $D^j f$ denotes the $j$-th derivative in the relevant Gaussian or Cameron--Martin directions, one asks whether
\[
\|f\|_{L^p}+\sum_{j=1}^{m}\|D^j f\|_{L^p}
\]
is equivalent to
\[
\|f\|_{L^p}+\|D^m f\|_{L^p}.
\]
The present paper is devoted to the second problem.  Approximation and density connect the two questions, but they are logically separate.

Sobolev analysis for Gaussian measures has its origins in the work of Gross on abstract Wiener spaces and infinite-dimensional potential theory \cite{GrossAWS,GrossPotential}.  Related elliptic and parabolic problems in infinitely many variables were studied by Daletskii \cite{Daletskii1967}; coercive estimates were obtained by Frolov \cite{Frolov1973}; and Kr\'ee developed Sobolev-type spaces and trace theory in infinite dimensions \cite{Kree1977}.  Gross's logarithmic Sobolev inequality \cite{Gross} subsequently made clear the importance of estimates whose constants are independent of the ambient dimension.  That feature is essential here: an inequality on each $\rr^n$ is useful for an infinite-dimensional limit only if the constants remain controlled as $n\to\infty$.  A systematic theory of Gaussian Sobolev spaces on locally convex spaces was later developed by Feyel and de La Pradelle \cite{FeyelPradelle1989}; further background can be found in \cite{BogachevGaussian,BogachevSurvey,DP}.

A parallel development came from Malliavin calculus.  Malliavin introduced the stochastic calculus of variations in his probabilistic approach to H\"ormander's theorem \cite{Malliavin1978}; analytic foundations were developed, among others, by Shigekawa \cite{Shigekawa1980}, Stroock \cite{Stroock1981}, and Watanabe \cite{Watanabe1984}.  In this setting the Malliavin derivative plays the role of the gradient, and the spaces $\mathbb D^{m,p}$ are the natural Gaussian Sobolev spaces.

For $1<p<\infty$, norm equivalence in Malliavin spaces is closely tied to the Ornstein--Uhlenbeck operator.  Meyer's work on Gaussian Riesz transforms \cite{Meyer1982,Meyer1983,Meyer1984} implies, by iteration, control of the intermediate Malliavin derivatives by the zeroth- and highest-order terms.  In the scalar case this yields the equivalence of the full Malliavin--Sobolev norm and the graph norm of the highest derivative.  Sugita extended the theory to Hilbert-space-valued Wiener functionals and clarified the structure of Sobolev spaces on abstract Wiener spaces \cite{Sugita1985,SugitaChar1985}.  Pisier gave a short analytic proof of Meyer's Riesz-transform inequality \cite{Pisier1988}, while Gundy developed a probabilistic treatment \cite{Gundy1989}.  Malliavin calculus and the corresponding estimates have also been extended to UMD Banach-space-valued random variables; see Pronk and Veraar \cite{PronkVeraar}.

The endpoint $p=1$ does not follow from this theory.  Meyer's inequalities are genuinely $L^p$ estimates in the range $1<p<\infty$, so the classical argument cannot simply be continued to $p=1$.

A second obstruction appears on proper subsets of an infinite-dimensional Gaussian space.  For convex sets in Wiener spaces, Hino studied the associated Dirichlet spaces and proved density results for smooth cylindrical functions under natural $H$-convexity and $H$-openness assumptions \cite{Hino2003,Hino2011}.  In finite dimensions, norm estimates on regular domains are often reduced to whole-space estimates by bounded extension operators.  Such a reduction is unavailable in general in infinite dimensions: Bogachev, Pilipenko, and Shaposhnikov showed, in particular, that Sobolev extension from convex Gaussian domains to the whole ambient space can fail \cite{BogachevPilipenkoShaposhnikov}.

There is by now a substantial regularity theory for differential equations on infinite-dimensional domains.  Da Prato and Lunardi established second-order and maximal Sobolev regularity for important elliptic and Neumann problems \cite{DaPratoLunardi2014,DaPratoLunardi2015}.  Cappa and Ferrari treated maximal regularity for weighted Gaussian measures, including convex subsets \cite{CappaFerrari2016,CappaFerrari2018}; related gradient estimates for perturbed Ornstein--Uhlenbeck semigroups on convex domains were obtained by Angiuli, Ferrari, and Pallara \cite{AngiuliFerrariPallara2019}, and weighted Gaussian Sobolev spaces were studied further by Ferrari \cite{Ferrari2019}.  These results concern density, traces, boundary conditions, operator domains, Dirichlet forms, and elliptic regularity.  They do not in general imply the higher-order norm equivalence considered here on an arbitrary proper open convex subset.

The endpoint problem for Malliavin--Sobolev norms was studied directly by Addona, Muratori, and Rossi \cite{AddonaMuratoriRossi}.  If $\mathfrak H$ is a real separable Hilbert space and $W$ is an isonormal Gaussian process over $\mathfrak H$, they compare
\[
\|F\|_{\mathcal D(k,p)}
=
\|F\|_{L^p}
+
\sum_{j=1}^{k}
\|D^jF\|_{L^p(\mathfrak H^{\otimes j})}
\]
with
\[
\|F\|_{\mathcal G(k,p)}
=
\|F\|_{L^p}
+
\|D^kF\|_{L^p(\mathfrak H^{\otimes k})}.
\]
For $1<p<\infty$ they recovered the known infinite-dimensional equivalence with quantitative control of the constants.  At $p=1$ they proved the case $k=2$ by a vector-valued Poincar\'e inequality.  The range
\[
p=1,\qquad k\geqslant3,
\]
was left open in \cite{AddonaMuratoriRossi}.  Their finite-dimensional argument applies to every $p\in[1,\infty)$, but the resulting constants grow with the dimension and therefore do not pass directly to a general infinite-dimensional Gaussian space.

More recently, Addona proved a whole-space result for weighted Gaussian measures $d\nu=K e^{-U}\,d\mu$ on separable Banach spaces, for $1<p<\infty$ and arbitrary differentiation order \cite{AddonaWeighted2025}.  The result substantially enlarges the class of admissible measures, but it remains a whole-space theorem and does not include the endpoint $p=1$.  The complementary ``$H=W$'' problem on subsets of $\ell^2$ is treated in \cite{WYZZ}, where density of suitable smooth functions in weak Sobolev spaces is established.  That density result will be used below to pass from cylindrical functions to Sobolev functions on a domain.

The purpose of this paper is to obtain finite-dimensional higher-order estimates with constants that are uniform in the dimension and then use them in two infinite-dimensional settings.

Let $\Omega_n\subset\rr^n$ be a nonempty open convex set, let $\gamma_n$ denote the standard Gaussian measure, and, for $j\geqslant0$, set
\[
A_j(f)
=
\left(
\int_{\Omega_n}\|D^jf\|_j^p\,d\gamma_n
\right)^{1/p},
\qquad D^0f=f.
\]
For every integer $m\geqslant2$ we prove an estimate of the form
\[
\sum_{j=1}^{m-1}A_j(f)
\leqslant
C\bigl(A_0(f)+A_m(f)\bigr),
\]
where the dependence of $C$ on the domain enters through $\gamma_n(\Omega_n)$, and there is no further dependence on $n$.

The proof has two components.  First, a vector-valued Gaussian Poincar\'e inequality is transferred to an arbitrary open convex set by Caffarelli's contraction theorem; see, for example, \cite{GozlanJuillet}.  Iterating this estimate for $D^rf,D^{r-1}f,\ldots,Df,f$ produces a polynomial $P$ of degree at most $r$ such that every derivative of $f-P$ of order at most $r$ is controlled by $D^{r+1}f$.  The remaining task is therefore a derivative estimate for a finite-degree polynomial.

Second, on the whole Gaussian space we estimate derivatives of polynomials by combining the Hermite expansion with Gaussian hypercontractivity.  The range $p\geqslant2$ follows directly.  For $1\leqslant p<2$, the $L^2$ estimate is combined with the Paley--Zygmund inequality.  The Carbery--Wright distributional inequality \cite{CW} then compares the global Gaussian $L^p$ norm of the polynomial with its $L^p$ norm on $\Omega_n$.  The resulting constants depend on the degree, on $p$, and on $\gamma_n(\Omega_n)$, but not otherwise on $n$.

The first infinite-dimensional application concerns a nondegenerate centered Gaussian measure on $\ell^2$.  Fix
\[
a_i>0,\qquad \sum_{i=1}^{\infty}a_i^2<\infty,
\]
and let $P$ be the Gaussian measure whose $i$-th coordinate has variance $a_i^2$.  If $\Omega\subset\ell^2$ is nonempty, open, and convex, define
\[
\Omega_n
=
\left\{
(x_1,\ldots,x_n)\in\rr^n:
(a_1x_1,\ldots,a_nx_n,\mathbf0^n)\in\Omega
\right\}.
\]
A key point is
\[
\gamma_n(\Omega_n)\longrightarrow P(\Omega)>0.
\]
Thus the domain-dependent constants from the finite-dimensional theorem remain uniformly bounded along the approximation.  Using cylindrical density, we obtain, for $m\geqslant2$,
\[
\begin{aligned}
\|f\|_{W^{m,p}(\Omega,P)}
\leqslant C_{\Omega,p,m}
\Bigg[
&\left(\int_\Omega |f|^p\,dP\right)^{1/p}
\\
&+
\left(
\int_\Omega
\left(
\sum_{i_1,\ldots,i_m=1}^{\infty}
 a_{i_1}^2\cdots a_{i_m}^2
|\partial_{i_1\cdots i_m}f|^2
\right)^{p/2}
 dP
\right)^{1/p}
\Bigg].
\end{aligned}
\]
The reverse inequality is immediate from the definition of the full Sobolev norm.  This gives the desired graph-norm characterization on a proper open convex subset of $\ell^2$.

The same finite-dimensional estimate also applies to Malliavin spaces.  When $\Omega_n=\rr^n$, one has $\gamma_n(\Omega_n)=1$, so the interpolation constants depend only on $p$ and the differentiation orders.  A smooth cylindrical random variable is a smooth function of finitely many independent standard Gaussian variables; the dimension-free estimate therefore passes directly to an arbitrary isonormal Gaussian process.  For every $p\in[1,\infty)$ and every $k\geqslant2$ we obtain
\[
\|F\|_{\mathcal G(k,p)}
\leqslant
\|F\|_{\mathcal D(k,p)}
\leqslant
C_{p,k}\|F\|_{\mathcal G(k,p)}.
\]
For $1<p<\infty$ this gives another proof of the classical equivalence associated with Meyer's inequalities.  For $p=1$, $k=2$, it recovers the endpoint result of Addona--Muratori--Rossi; for $p=1$, $k\geqslant3$, it settles the case left open in \cite{AddonaMuratoriRossi}.

We also prove the Hilbert-space-valued version.  A coordinatewise application of the scalar estimate would introduce a dependence on the dimension of the range.  Gaussian randomization and the Gaussian Kahane--Khintchine inequality \cite{HNVW} avoid this loss and give
\[
\|F\|_{\mathcal G(k,p)(V)}
\leqslant
\|F\|_{\mathcal D(k,p)(V)}
\leqslant
K_{p,k}\|F\|_{\mathcal G(k,p)(V)},
\]
with $K_{p,k}$ independent of both $\dim\mathfrak H$ and $\dim V$.  Consequently, the completions defined by the two norms coincide.

The argument brings together several dimension-free features of Gaussian analysis: Gaussian Poincar\'e inequalities, contraction for Gaussian transport, hypercontractivity, dimension-independent polynomial estimates such as those in \cite{EI}, and the Carbery--Wright inequality \cite{CW}.  In both applications the same principle is used: the infinite-dimensional conclusion is obtained only after the finite-dimensional constants have been made independent of the dimension.

The paper is organized as follows.  Section~2 fixes the notation and the Sobolev structures used later.  Section~3 proves the dimension-free estimates on open convex subsets of $\rr^n$: the vector-valued Poincar\'e inequality, the polynomial approximation, and the derivative estimate for finite-degree polynomials are established in that order.  Section~4 passes these estimates to open convex subsets of $\ell^2$.  Section~5 treats scalar- and Hilbert-space-valued Malliavin spaces and proves norm equivalence for all $p\in[1,\infty)$ and $k\geqslant2$.


\section{Preliminaries}

Throughout the paper, $\nn=\{1,2,\ldots\}$ and $\Nzero=\nn\cup\{0\}$.  If $X$ is a topological space, $\mathscr B(X)$ denotes its Borel $\sigma$-algebra.  For real Banach spaces $X$ and $Y$, a nonempty open set $O\subset X$, and $k\in\Nzero$, we write $C^k(O;Y)$ for the space of $k$-times continuously Fr\'echet differentiable maps from $O$ to $Y$.  In the scalar case, $C^k(O)$ means $C^k(O;\rr)$.

Let $C_P^\infty(\rr^n)$ be the collection of all $C^\infty$ functions on $\rr^n$ whose derivatives of every order have at most polynomial growth.  Identifying a function on $\rr^n$ with the corresponding cylindrical function on $\ell^2$, set
\[
\mathscr C_P^\infty
\triangleq
\bigcup_{n=1}^\infty C_P^\infty(\rr^n).
\]
For a set $A$, $\chi_A$ denotes its indicator function.  Unless stated otherwise, $p\in[1,\infty)$.

Fix a nonempty open convex set $\Omega_n\subset\rr^n$.  Let $\mathbf e_1,\ldots,\mathbf e_n$ be the standard basis of $\rr^n$ and $\mathbf e_1^*,\ldots,\mathbf e_n^*$ the corresponding dual basis.  Put
\[
\Ls{0}(\rr^n)\triangleq\rr,
\qquad
\langle f,g\rangle_0\triangleq fg,
\qquad
\|f\|_0\triangleq |f|.
\]
For $k\in\nn$, let $\Ls{k}(\rr^n)$ be the vector space of symmetric $k$-linear forms on $\rr^n$.  If
\[
f=
\sum_{i_1,\ldots,i_k=1}^n
c_{i_1\cdots i_k}\,
\mathbf e_{i_1}^*\otimes\cdots\otimes\mathbf e_{i_k}^*,
\]
and
\[
g=
\sum_{i_1,\ldots,i_k=1}^n
 d_{i_1\cdots i_k}\,
\mathbf e_{i_1}^*\otimes\cdots\otimes\mathbf e_{i_k}^*,
\]
define
\[
\langle f,g\rangle_k
\triangleq
\sum_{i_1,\ldots,i_k=1}^n
c_{i_1\cdots i_k}d_{i_1\cdots i_k},
\qquad
\|f\|_k
\triangleq
\left(
\sum_{i_1,\ldots,i_k=1}^n c_{i_1\cdots i_k}^2
\right)^{1/2}.
\]
Thus $\Ls{k}(\rr^n)$ is a finite-dimensional Hilbert space.  We write $\mathcal P(\rr^n)$ for the real polynomials on $\rr^n$.

For $k\in\nn$ and $(i_1,\ldots,i_k)\in\{1,\ldots,n\}^k$, we use
\[
\partial_{i_1\cdots i_k}f
\triangleq
\frac{\partial^k f}{\partial x_{i_1}\cdots\partial x_{i_k}}.
\]
The standard Gaussian measure on $\rr^n$ is
\[
d\gamma_n(\mathbf x)
\triangleq
(2\pi)^{-n/2}e^{-|\mathbf x|^2/2}\,d\mathbf x.
\]
For a sufficiently smooth scalar function $f$, define
\[
A_0(f)
\triangleq
\left(\int_{\Omega_n}|f|^p\,d\gamma_n\right)^{1/p},
\]
and, for $k\in\nn$,
\[
A_k(f)
\triangleq
\left[
\int_{\Omega_n}
\left(
\sum_{i_1,\ldots,i_k=1}^n
|\partial_{i_1\cdots i_k}f|^2
\right)^{p/2}
 d\gamma_n
\right]^{1/p}.
\]
Equivalently, $A_k(f)=\|\,\|D^kf\|_k\,\|_{L^p(\Omega_n,\gamma_n)}$.

\section{Dimension-free Sobolev estimates on convex Gaussian domains}
Throughout this section, $D^0f=f$, and for $k\in\nn$ the symbol $D^kf$ denotes the $k$-th Fr\'echet derivative.  Its norm is the Hilbert--Schmidt norm induced by $\langle\cdot,\cdot\rangle_k$.
\begin{lemma}\label{qi}
Let $k\in\nn$ and $f\in C^{k+1}(\rr^n)$. Then
\begin{eqnarray}
\|D^{k+1}f\|_{k+1}
=
\left(\sum_{i=1}^{n}\|\partial_iD^kf\|_k^2\right)^{1/2}.
\label{qiudao}
\end{eqnarray}
\end{lemma}

\begin{remark}
Here $\partial_i(D^kf)$ denotes $D(D^kf)\mathbf e_i$, regarded as an element of $\Ls{k}(\rr^n)$.
\end{remark}

\begin{proof}
Writing
\[
D^kf
=
\sum_{i_1,\ldots,i_k=1}^n
\partial_{i_1\cdots i_k}f\,
\mathbf e_{i_1}^*\otimes\cdots\otimes\mathbf e_{i_k}^*,
\]
we obtain
\[
\begin{aligned}
\sum_{i=1}^n\|\partial_iD^kf\|_k^2
&=
\sum_{i=1}^n\sum_{i_1,\ldots,i_k=1}^n
\left|\partial_i\partial_{i_1\cdots i_k}f\right|^2
\\
&=
\|D^{k+1}f\|_{k+1}^2.
\end{aligned}
\]
Taking square roots proves Lemma~\ref{qi}.
\end{proof}

The next result transfers the vector-valued Gaussian Poincar\'e inequality to an open convex Gaussian domain.  Its constant is independent of the ambient dimension.

\begin{theorem}\label{poincare}
Let $(H,\langle\cdot,\cdot\rangle_H)$ be a finite-dimensional real Hilbert space.  Suppose that $F\in C^\infty(\Omega_n;H)$ and that $\|F\|_H$ and $\|\partial_iF\|_H$, $i=1,\ldots,n$, have at most polynomial growth.  Then
\begin{eqnarray}
\int_{\Omega_n}
\left\|
F-
\frac{1}{\gamma_n(\Omega_n)}
\int_{\Omega_n}F\,d\gamma_n
\right\|_H^p
\,d\gamma_n
\leqslant
k_p
\int_{\Omega_n}
\left(
\sum_{i=1}^n\|\partial_iF\|_H^2
\right)^{p/2}
\,d\gamma_n,
\label{1}
\end{eqnarray}
where
\begin{eqnarray}
k_p
\triangleq
\begin{cases}
(p-1)^{p/2},&p\in[2,\infty),\\[2mm]
(\pi/2)^p,&p\in[1,2).
\end{cases}
\label{2}
\end{eqnarray}
\end{theorem}

\begin{proof}
Let
\[
d\mu_n
\triangleq
\frac{\chi_{\Omega_n}}{\gamma_n(\Omega_n)}\,d\gamma_n.
\]
Equivalently,
\[
d\mu_n=e^{-V_n}\,d\gamma_n,
\qquad
V_n=-\log\chi_{\Omega_n}+\log\gamma_n(\Omega_n).
\]
Because $\Omega_n$ is convex, $V_n:\rr^n\to\rr\cup\{+\infty\}$ is convex.  Caffarelli's contraction theorem, applied in the extended-valued case by the usual convex approximation, yields a $1$-Lipschitz transport map $T:\rr^n\to\rr^n$ such that
\[
T_{\#}\gamma_n=\mu_n;
\]
see, for instance, \cite[Theorem~2.2, p.~440]{GozlanJuillet}.

Apply the vector-valued Gaussian Poincar\'e inequality \cite[Theorem~2.6, p.~9]{AddonaMuratoriRossi} to $F\circ T$.  The composition is interpreted $\gamma_n$-almost everywhere; since $T_{\#}\gamma_n=\mu_n$ and $\mu_n$ is concentrated on $\Omega_n$, this is sufficient.  A standard Sobolev approximation gives
\[
\begin{aligned}
&\int_{\rr^n}
\left\|
F\circ T-
\int_{\rr^n}F\circ T\,d\gamma_n
\right\|_H^p
\,d\gamma_n
\\
&\qquad\leqslant
k_p
\int_{\rr^n}
\left(
\sum_{i=1}^n
\|\partial_i(F\circ T)\|_H^2
\right)^{p/2}
\,d\gamma_n.
\end{aligned}
\]
Since $T$ is $1$-Lipschitz, Rademacher's theorem gives $\|DT\|_{\mathrm{op}}\leqslant1$ almost everywhere.  The Sobolev chain rule therefore yields
\[
\left(
\sum_{i=1}^n\|\partial_i(F\circ T)\|_H^2
\right)^{1/2}
\leqslant
\left(
\sum_{i=1}^n\|(\partial_iF)\circ T\|_H^2
\right)^{1/2}
\quad\text{a.e.}
\]
Consequently,
\[
\begin{aligned}
&\int_{\rr^n}
\left\|
F\circ T-
\int_{\rr^n}F\circ T\,d\gamma_n
\right\|_H^p
\,d\gamma_n
\\
&\qquad\leqslant
k_p
\int_{\rr^n}
\left(
\sum_{i=1}^n
\|(\partial_iF)\circ T\|_H^2
\right)^{p/2}
\,d\gamma_n.
\end{aligned}
\]
Using $T_{\#}\gamma_n=\mu_n$ on both sides gives
\[
\int_{\Omega_n}
\left\|
F-\int_{\Omega_n}F\,d\mu_n
\right\|_H^p
\,d\mu_n
\leqslant
k_p
\int_{\Omega_n}
\left(
\sum_{i=1}^n\|\partial_iF\|_H^2
\right)^{p/2}
\,d\mu_n.
\]
Multiplying by $\gamma_n(\Omega_n)$ and using the definition of $\mu_n$ gives \eqref{1}.  This proves Theorem~\ref{poincare}.
\end{proof}

\begin{theorem}\label{bijin}
Let $r\in\mathbb{N}_0$ and $f\in C_P^{\infty}(\mathbb{R}^{n})$. Then there exists $P\in \mathcal{P} \left( \mathbb{R}^{n} \right)$ such that $\deg P\leqslant r$ and, for $j=0,1,2,\cdots ,r$,
\begin{eqnarray}
	A_{j}\left( f-P \right) \leqslant  k_{p}^{\frac{r+1-j}{p}}A_{r+1}\left( f \right).\label{233}
\end{eqnarray}
\end{theorem}
\begin{proof}
If $r=0$, take
$$T_0\triangleq\frac{1}{\gamma_n(\Omega_n)}\int_{\Omega_n}f\,\mathrm d\gamma_n,
\qquad P\triangleq T_0.$$
Then \eqref{233} is exactly Theorem~\ref{poincare} applied to the scalar function $f$.  Hence we may assume $r\geqslant1$.

We construct the coefficients of the polynomial recursively.  First, define
$$T_{r}\triangleq \frac{1}{\gamma_{n} \left( \Omega_n \right)} \int_{\Omega_n} D^{r}f\mathrm{d} \gamma_{n} \in \mathcal{L}_{s}^{r} \left( \mathbb{R}^{n} \right).$$
Suppose next that $T_r,T_{r-1},\ldots,T_k$ have already been defined for some $1\leqslant k\leqslant r$.  Define
$$T_{k-1}\triangleq \frac{1}{\gamma_{n} \left( \Omega_n \right)} \int_{\Omega_n} D^{k-1}\left( f-\frac{T_{r}\left( \textbf{x} ,\cdots ,\textbf{x} \right)}{r!} -\cdots -\frac{T_{k}\left( \textbf{x} ,\cdots ,\textbf{x} \right)}{k!} \right) \mathrm{d} \gamma_{n} \in \mathcal{L}_{s}^{k-1} \left( \mathbb{R}^{n} \right).$$
For $k=1,2,\cdots ,r$, it follows from Lemma \ref{qi} and Theorem \ref{poincare} that$$\begin{aligned}
	&\int_{\Omega_n} \left| \left| D^{k-1}\left( f-\frac{T_{r}\left( \textbf{x} ,\textbf{x} ,\cdots ,\textbf{x} \right)}{r!} -\cdots -\frac{T_{k}\left( \textbf{x} ,\cdots ,\textbf{x} \right)}{k!} \right) -T_{k-1} \right| \right|_{k-1}^{p} \mathrm{d} \gamma_{n}\\ &\  \leqslant k_{p}\int_{\Omega_n} \left( \sqrt{\sum_{i=1}^{n} \left| \left| \partial_{i} \left[ D^{k-1}\left( f-\frac{T_{r}\left( \textbf{x} ,\textbf{x} ,\cdots ,\textbf{x} \right)}{r!} -\cdots -\frac{T_{k}\left( \textbf{x} ,\cdots ,\textbf{x} \right)}{k!} \right) \right] \right| \right|_{k-1}^{2}} \right)^{p} \mathrm{d} \gamma_{n}\\ &\  =k_{p}\int_{\Omega_n} \left| \left| D^{k}\left( f-\frac{T_{r}\left( \textbf{x} ,\textbf{x} ,\cdots ,\textbf{x} \right)}{r!} -\cdots -\frac{T_{k}\left( \textbf{x} ,\cdots ,\textbf{x} \right)}{k!} \right) \right| \right|_{k}^{p} \mathrm{d} \gamma_{n}\\ &\  =k_{p}\int_{\Omega_n} \left| \left| D^{k}\left( f-\frac{T_{r}\left( \textbf{x} ,\textbf{x} ,\cdots ,\textbf{x} \right)}{r!} -\cdots -\frac{T_{k+1}\left( \textbf{x} ,\cdots ,\textbf{x} \right)}{\left( k+1 \right) !} \right) -T_{k} \right| \right|_{k}^{p} \mathrm{d} \gamma_{n} .\end{aligned}$$	Let $P\triangleq T_{0}+\sum_{j=1}^{r} \frac{T_{j}\left( \textbf{x} ,\textbf{x}, \cdots ,\textbf{x} \right)}{j!}\in \mathcal{P} \left( \mathbb{R}^{n} \right)$.  Then $\deg P\leqslant r$.  For $j=0,1,2,\ldots,r$, iterating the preceding inequality gives
$$\begin{aligned}
	&A_{j}\left( f-P \right) =\sqrt[p]{\int_{\Omega_{n}} \left| \left| D^{j}\left( f-P \right) \right| \right|_{j}^{p} \mathrm{d} \gamma_{n}}\\ &\  =\sqrt[p]{\int_{\Omega_{n}} \left\vert \left\vert D^{j}\left( f-\sum_{k=j}^{r} \frac{T_{k}\left( \textbf{x} ,\textbf{x} ,\cdots ,\textbf{x} \right)}{k!} \right) \right\vert \right\vert_{j}^{p} \mathrm{d} \gamma_{n}}\\ &\  \leqslant \sqrt[p]{k_{p}} \sqrt[p]{\int_{\Omega_{n}} \left\vert \left\vert D^{j+1}\left( f-\sum_{k=j+1}^{r} \frac{T_{k}\left( \textbf{x} ,\textbf{x} ,\cdots ,\textbf{x} \right)}{k!} \right) \right\vert \right\vert_{j+1}^{p} \mathrm{d} \gamma_{n}}\\ &\  \  \  \vdots\\ &\  \leqslant \left( \sqrt[p]{k_{p}} \right)^{r-j} \sqrt[p]{\int_{\Omega_{n}} \left\vert \left\vert D^{r}\left( f-\frac{T_{r}\left( \textbf{x} ,\textbf{x} ,\cdots ,\textbf{x} \right)}{r!} \right) \right\vert \right\vert_{r}^{p} \mathrm{d} \gamma_{n}}\\ &\  =\left( \sqrt[p]{k_{p}} \right)^{r-j} \sqrt[p]{\int_{\Omega_{n}} \left\vert \left\vert D^{r}f-T_{r} \right\vert \right\vert_{r}^{p} \mathrm{d} \gamma_{n}}\\ &\  \leqslant \left( \sqrt[p]{k_{p}} \right)^{r+1-j} A_{r+1}\left( f \right) ,\end{aligned}$$
This proves Theorem~\ref{bijin}.
\end{proof}
\begin{theorem}\label{weishuguanguji}
	Let $r,k\in \mathbb{N}$. Then there exists $C_{p,r,k}>0$ such that, for every $P\in \mathcal{P} \left( \mathbb{R}^{n} \right)$ of degree at most $r$, we have
	\begin{eqnarray}
	\int_{\Omega_{n}} \left| \left| D^{k}P \right| \right|_{k}^{p} \mathrm{d} \gamma_{n} \leqslant \frac{C_{p,r,k}}{[\gamma_n(\Omega_n)]^{1+pr}} \int_{\Omega_{n}} \left| P \right|^{p} \mathrm{d} \gamma_{n}.\label{1fff13f}
	\end{eqnarray}
\end{theorem}
\begin{proof}
If $P\equiv0$, the estimate is immediate.  We may therefore assume $P\not\equiv0$.  If $k>r$, then $D^kP\equiv0$, so the estimate again holds (for instance with $C_{p,r,k}=1$).  It remains to consider $1\leqslant k\leqslant r$.

\medskip\noindent\textit{Step 1: Hermite expansion and the $L^2$ derivative estimate.}

By \cite[Proposition~5.48, p.~135]{AS17},
\begin{eqnarray}
	\int_{\mathbb{R}^{n}} \left| P \right|^{q} \mathrm{d} \gamma_{n} \leqslant \left( q-1 \right)^{\frac{rq}{2}} \left( \int_{\mathbb{R}^{n}} \left| P \right|^{2} \mathrm{d} \gamma_{n} \right)^{\frac{q}{2}},\forall q\geqslant 2.\label{13f1f3}
\end{eqnarray}
	
	Let
	$$H_{m}\left( t \right) \triangleq \frac{\left( -1 \right)^{m}}{\sqrt{m!}} e^{\frac{t^{2}}{2}}\frac{d^{m}}{dt^{m}} e^{-\frac{t^{2}}{2}},m=0,1,2,\cdots,$$
	and
	$$H_{\alpha}\left( x_{1},x_{2},\cdots ,x_{n} \right) =\prod_{j=1}^{n} H_{\alpha_{j}}\left( x_{j} \right) ,\alpha = \left( \alpha_{j} \right) \in \mathbb{N}_{0}^{n}.$$
	We know that $\left\{ H_{\alpha} \right\}_{\alpha \in \mathbb{N}_{0}^{n}}$ is a complete orthonormal basis of $L^{2}\left( \mathbb{R}^{n} ,\gamma_{n} \right)$ and that
	\begin{eqnarray}
		\partial_{i} H_{\alpha}=\sqrt{\alpha_{i}} H_{\alpha -\textbf{e}_{i}},i=1,2,\cdots ,n,\alpha \in \mathbb{N}_{0}^{n}.\label{1f1vvv}
	\end{eqnarray}
	Consider $\left| \alpha \right| \triangleq \sum\limits_{i=1}^{n} \alpha_{i}$ and the expansion
	$$P=\sum_{\left| \alpha \right| \leqslant r} c_{\alpha}H_{\alpha},$$
	where $c_{\alpha}\in \mathbb{R} ,\forall \alpha \in \mathbb{N}_{0}^{n}$. We have
	$$\int_{\mathbb{R}^{n}} \left| P \right|^{2} \mathrm{d} \gamma_{n} =\sum_{\left| \alpha \right| \leqslant r} \left| c_{\alpha} \right|^{2}.$$
	For $(i_1,i_2,\ldots,i_k)\in\{1,2,\ldots,n\}^k,k=1,2,\cdots,r$, let $\beta_{j}$ denote the number of occurrences of $j$ among $i_{1},i_{2},\cdots ,i_{k}$. Using \eqref{1f1vvv}, we obtain
	$$\int_{\mathbb{R}^{n}} \left| \partial_{i_{1}i_{2}\cdots i_{k}} P \right|^{2} \mathrm{d} \gamma_{n} =\sum_{\left| \alpha \right| \leqslant r} \left| c_{\alpha} \right|^{2} \cdot \left[ \prod_{j=1}^{n} \alpha_{j} \left( \alpha_{j} -1 \right) \cdots \left( \alpha_{j} -\beta_{j} +1 \right) \right],$$
	where, when $\beta_{j} =0$, the corresponding product is understood to be $1$. 
	
	Note the identity
	\[
	\sum_{i_{1},i_{2},\ldots,i_{k}=1}^{n}
	\prod_{j=1}^{n}
	\alpha_{j}(\alpha_{j}-1)\cdots(\alpha_{j}-\beta_{j}+1)
	=
	|\alpha|(|\alpha|-1)\cdots(|\alpha|-k+1).
	\]
	Indeed, regard the $j$-th box as containing $\alpha_j$ distinct balls.
	For a fixed ordered sequence
	$(i_1,i_2,\ldots,i_k)\in\{1,2,\ldots,n\}^k$,
	the quantity
	\[
	\prod_{j=1}^{n}
	\alpha_{j}(\alpha_{j}-1)\cdots(\alpha_{j}-\beta_{j}+1)
	\]
	is the number of ordered selections of $k$ distinct balls such that the
	$\ell$-th selected ball is taken from the $i_\ell$-th box.
	Therefore, summing over all ordered sequences
	$(i_1,i_2,\ldots,i_k)$ counts each ordered selection of $k$ distinct balls
	from the total of $|\alpha|$ balls exactly once, and the total number of such
	selections is
	\[
	|\alpha|(|\alpha|-1)\cdots(|\alpha|-k+1).
	\]

We now have
	$$\begin{aligned}\sum_{i_{1},i_{2},\cdots, i_{k}=1}^{n} \int_{\mathbb{R}^{n}} \left| \partial_{i_{1}i_{2}\cdots i_{k}} P \right|^{2} \mathrm{d} \gamma_{n}&=\sum_{\left| \alpha \right| \leqslant r} \left[ \left| c_{\alpha} \right|^{2} \left| \alpha \right| \cdot \left( \left| \alpha \right| -1 \right) \cdots \left( \left| \alpha \right| -k+1 \right) \right]\\ &\leqslant r\left( r-1 \right) \cdots \left( r-k+1 \right) \sum_{\left| \alpha \right| \leqslant r} \left| c_{\alpha} \right|^{2}\\ &=r\left( r-1 \right) \cdots \left( r-k+1 \right) \int_{\mathbb{R}^{n}} \left| P \right|^{2} \mathrm{d} \gamma_{n}.\end{aligned}$$
	
	Hence, for $k=1,2,\cdots ,r$, we obtain the following basic estimate:
	\begin{eqnarray}
		\int_{\mathbb{R}^{n}} \left| \left| D^{k}P \right| \right|_{k}^{2} \mathrm{d} \gamma_{n} \leqslant r\left( r-1 \right) \cdots \left( r-k+1 \right) \int_{\mathbb{R}^{n}} \left| P \right|^{2} \mathrm{d} \gamma_{n}.\label{1f1f3}
	\end{eqnarray}
	
\medskip\noindent\textit{Step 2: Passage from $L^2$ to $L^p$.}

When $p\geqslant 2,k=1,2,\cdots,r$, we have
$$\begin{aligned}\left( \int_{\mathbb{R}^{n}} \left| \left| D^{k}P \right| \right|_{k}^{p} \mathrm{d} \gamma_{n} \right)^{\frac{2}{p}}&=\left( \int_{\mathbb{R}^{n}} \left( \sum_{i_{1},i_{2},\cdots ,i_{k}=1}^{n} \left| \partial_{i_{1}i_{2}\cdots i_{k}} P \right|^{2} \right)^{\frac{p}{2}} \mathrm{d} \gamma_{n} \right)^{\frac{2}{p}}\\ &\leqslant \sum_{i_{1},i_{2},\cdots ,i_{k}=1}^{n} \left( \int_{\mathbb{R}^{n}} \left( \left| \partial_{i_{1}i_{2}\cdots i_{k}} P \right|^{2} \right)^{\frac{p}{2}} \mathrm{d} \gamma_{n} \right)^{\frac{2}{p}}\\ &=\sum_{i_{1},i_{2},\cdots ,i_{k}=1}^{n} \left( \int_{\mathbb{R}^{n}} \left| \partial_{i_{1}i_{2}\cdots i_{k}} P \right|^{p} \mathrm{d} \gamma_{n} \right)^{\frac{2}{p}}\\ &\leqslant \sum_{i_{1},i_{2},\cdots ,i_{k}=1}^{n} \left( p-1 \right)^{r-k} \int_{\mathbb{R}^{n}} \left| \partial_{i_{1}i_{2}\cdots i_{k}} P \right|^{2} \mathrm{d} \gamma_{n}\\ &=\left( p-1 \right)^{r-k} \int_{\mathbb{R}^{n}} \left| \left| D^{k}P \right| \right|_{k}^{2} \mathrm{d} \gamma_{n}\\ &\leqslant \left( p-1 \right)^{r-k} r\left( r-1 \right) \cdots \left( r-k+1 \right) \int_{\mathbb{R}^{n}} \left| P \right|^{2} \mathrm{d} \gamma_{n}\\ &\leqslant \left( p-1 \right)^{r-k} r\left( r-1 \right) \cdots \left( r-k+1 \right) \left( \int_{\mathbb{R}^{n}} \left| P \right|^{p} \mathrm{d} \gamma_{n} \right)^{\frac{2}{p}} ,\end{aligned}$$
where the first inequality follows from $\frac{p}{2} \geqslant 1$ and Minkowski's inequality, the second inequality follows from \eqref{13f1f3}, the third inequality follows from \eqref{1f1f3}, and the last inequality follows from Hölder's inequality. Therefore, for $p\in \left[ 2,+\infty \right),k=1,2,\cdots ,r$, we obtain the estimate
\begin{eqnarray}
	\int_{\mathbb{R}^{n}} \left| \left| D^{k}P \right| \right|_{k}^{p} \mathrm{d} \gamma_{n} \leqslant \left( p-1 \right)^{\frac{r-k}{2} p} \left( r\left( r-1 \right) \cdots \left( r-k+1 \right) \right)^{\frac{p}{2}} \int_{\mathbb{R}^{n}} \left| P \right|^{p} \mathrm{d} \gamma_{n}.\label{11}
\end{eqnarray}	

When $1\leqslant p < 2,k=1,2,\cdots,r$, it follows first from \eqref{13f1f3} that
$$\int_{\mathbb{R}^{n}} \left\vert P \right\vert^{4} \mathrm{d} \gamma_{n} \leqslant 9^{r}\left( \int_{\mathbb{R}^{n}} \left| P \right|^{2} \mathrm{d} \gamma_{n} \right)^{2}.$$
Then, by the Paley--Zygmund inequality (see, for instance, \cite[Proposition~III.3, p.~23]{LQ17}), we have
\begin{eqnarray}\label{1ff133}
	\begin{aligned}
		&\gamma_{n} \left( \left\{ \textbf{x} \in \mathbb{R}^{n} :\left| P \right| \geqslant \sqrt{\frac{1}{2} \int_{\mathbb{R}^{n}} \left| P \right|^{2} \mathrm{d} \gamma_{n}} \right\} \right)\\ &\  \geqslant \frac{\left( \int_{\mathbb{R}^{n}} \left| P \right|^{2} \mathrm{d} \gamma_{n} \right)^{2}}{4\int_{\mathbb{R}^{n}} \left\vert P \right\vert^{4} \mathrm{d} \gamma_{n}} \geqslant \frac{1}{4\cdot 9^{r}}.\end{aligned}
\end{eqnarray}
Hence, together with Hölder's inequality,
$$\left( \int_{\mathbb{R}^{n}} \left| P \right|^{p} \mathrm{d} \gamma_{n} \right)^{\frac{1}{p}} \geqslant \int_{\mathbb{R}^{n}} \left| P \right| \mathrm{d} \gamma_{n} \geqslant \sqrt{\frac{1}{2} \int_{\mathbb{R}^{n}} \left| P \right|^{2} \mathrm{d} \gamma_{n}} \cdot \frac{1}{4\cdot 9^{r}}.$$
Therefore, we have
$$\begin{aligned}
	&\int_{\mathbb{R}^{n}} \left| \left| D^{k}P \right| \right|_{k}^{p} \mathrm{d} \gamma_{n} \leqslant \left( \int_{\mathbb{R}^{n}} \left| \left| D^{k}P \right| \right|_{k}^{2} \mathrm{d} \gamma_{n} \right)^{\frac{p}{2}}\\ &\  \leqslant \left( r\left( r-1 \right) \cdots \left( r-k+1 \right) \right)^{\frac{p}{2}} \left( \int_{\mathbb{R}^{n}} \left\vert P \right\vert^{2} \mathrm{d} \gamma_{n} \right)^{\frac{p}{2}}\\ &\  \leqslant \left( 4\sqrt{2} \right)^{p} \cdot 3^{2pr}\cdot \left( r\left( r-1 \right) \cdots \left( r-k+1 \right) \right)^{\frac{p}{2}} \int_{\mathbb{R}^{n}} \left\vert P \right\vert^{p} \mathrm{d} \gamma_{n},\end{aligned}$$
where the first inequality follows from Hölder's inequality, the second inequality follows from \eqref{1f1f3}, and the third inequality follows from \eqref{1ff133}.

Therefore, for $p\in \left[ 1,2 \right),k=1,2,\cdots ,r$, we obtain the estimate
\begin{eqnarray}
	\int_{\mathbb{R}^{n}} \left| \left| D^{k}P \right| \right|_{k}^{p} \mathrm{d} \gamma_{n} \leqslant \left( 4\sqrt{2} \right)^{p} \cdot 3^{2pr}\cdot\left( r\left( r-1 \right) \cdots \left( r-k+1 \right) \right)^{\frac{p}{2}} \int_{\mathbb{R}^{n}} \left| P \right|^{p} \mathrm{d} \gamma_{n}.\label{12}
\end{eqnarray}
\medskip\noindent\textit{Step 3: Localization from $\rr^n$ to $\Omega_n$.}
	
The Carbery--Wright inequality \cite[Theorem~8, p.~244]{CW}, applied with degree bound $d=r$ and parameter $q=rp$, yields an absolute constant $C>0$ such that, for every $\alpha>0$,
$$\gamma_{n} \left\{ \textbf{x} \in \mathbb{R}^{n} :\left| P\left( \textbf{x} \right) \right| \leqslant \alpha \right\} \leqslant Crp\alpha^{\frac{1}{r}} \left( \int_{\mathbb{R}^{n}} \left| P \right|^{p} \mathrm{d} \gamma_{n} \right)^{-\frac{1}{rp}}.$$
Choosing $\alpha_0 =\left( \frac{\gamma_{{}n} \left( \Omega_{n} \right)}{2Crp} \right)^{r} \sqrt[p]{\int_{\mathbb{R}^{n}} \left| P \right|^{p} \mathrm{d} \gamma_{n}}$ and substituting it into the preceding inequality gives
\begin{eqnarray}
	\gamma_{n} \left\{ \textbf{x} \in \mathbb{R}^{n} :\left| P\left( \textbf{x} \right) \right| \leqslant \alpha_{0} \right\} \leqslant \frac{\gamma_{{}n} \left( \Omega_{n} \right)}{2}.\label{13f13ff}
\end{eqnarray}
Now, by \eqref{13f13ff},
$$\begin{aligned}
	&\gamma_{n} \left( \Omega_{n} \bigcap \left\{ \textbf{x} \in \mathbb{R}^{n} :\left| P\left( \textbf{x} \right) \right| >\alpha_{0} \right\} \right)\\ &\  =\gamma_{n} \left( \Omega_{n} \right) -\gamma_{n} \left( \Omega_{n} \bigcap \left\{ \textbf{x} \in \mathbb{R}^{n} :\left| P\left( \textbf{x} \right) \right| \leqslant \alpha_{0} \right\} \right)\\ &\  \geqslant \gamma_{n} \left( \Omega_{n} \right) -\frac{\gamma_{n} \left( \Omega_{n} \right)}{2} =\frac{\gamma_{n} \left( \Omega_{n} \right)}{2} .\end{aligned}$$
Therefore,
$$\begin{aligned}\int_{\Omega_{n}} \left| P \right|^{p} \mathrm{d} \gamma_{n}&\geqslant \int_{\Omega_{n} \bigcap \left\{ \textbf{x} \in \mathbb{R}^{n} :\left| P \right| >\alpha_{0} \right\}} \left| P \right|^{p} \mathrm{d} \gamma_{n} \geqslant \frac{\gamma_{{}n} \left( \Omega_{n} \right)}{2} \alpha_{0}^{p}\\ &=\frac{\gamma_{{}n} \left( \Omega_{n} \right)}{2} \left( \frac{\gamma_{{}n} \left( \Omega_{n} \right)}{2Crp} \right)^{pr} \int_{\mathbb{R}^{n}} \left| P \right|^{p} \mathrm{d} \gamma_{n} ,\end{aligned}$$
that is,
\begin{eqnarray}
\int_{\mathbb{R}^{n}} \left| P \right|^{p} \mathrm{d} \gamma_{n} \leqslant \left( \frac{2}{\gamma_{n} \left( \Omega_{n} \right)} \right) \left( \frac{2Crp}{\gamma_{n} \left( \Omega_{n} \right)} \right)^{pr} \int_{\Omega_{n}} \left| P \right|^{p} \mathrm{d} \gamma_{n}.\label{1f33gfg}
\end{eqnarray}

\medskip\noindent\textit{Step 4: Conclusion.}

Combining \eqref{11}, \eqref{12}, and \eqref{1f33gfg}, we obtain a constant $C_{p,r,k}>0$ such that \eqref{1fff13f} holds. This completes the proof of Theorem~\ref{weishuguanguji}.

\end{proof}
\begin{theorem}\label{zyao}
Let $r\in\mathbb{N}$ and $f\in C_P^{\infty}(\mathbb{R}^{n})$. Then
	\begin{eqnarray}
\begin{aligned}
\sum_{j=1}^{r}A_j(f)
&\leqslant
\sum_{j=1}^{r}
\max\left\{
\frac{C_{p,r,j}^{1/p}}{[\gamma_n(\Omega_n)]^{1/p+r}},
\frac{C_{p,r,j}^{1/p}}{[\gamma_n(\Omega_n)]^{1/p+r}}k_p^{(r+1)/p}
+k_p^{(r+1-j)/p}
\right\}\\
&\qquad\qquad\times\bigl(A_0(f)+A_{r+1}(f)\bigr).
\end{aligned}
\label{1f1f13f3f1}
\end{eqnarray}
\end{theorem}
\begin{proof}
Let $P$ be the polynomial furnished by Theorem~\ref{bijin}.  For each $j=1,2,\ldots,r$, Theorems~\ref{bijin} and~\ref{weishuguanguji} give
	$$\begin{aligned}A_{j}\left( f \right)&\leqslant A_{j}\left( f-P \right) +A_{j}\left( P \right) \leqslant k_{p}^{\frac{r+1-j}{p}}A_{r+1}\left( f \right) +\frac{C_{p,r,j}^{\frac{1}{p}}}{[\gamma_n(\Omega_n)]^{\frac{1}{p}+r}} \sqrt[p]{\int_{\Omega_{n}} \left| P \right|^{p} \mathrm{d} \gamma_{n}}\\ &\leqslant k_{p}^{\frac{r+1-j}{p}}A_{r+1}\left( f \right) +\frac{C_{p,r,j}^{\frac{1}{p}}}{[\gamma_n(\Omega_n)]^{\frac{1}{p}+r}} \left[ A_{0}\left( f \right) +A_{0}\left( f-P \right) \right]\\ &\leqslant k_{p}^{\frac{r+1-j}{p}}A_{r+1}\left( f \right) +\frac{C_{p,r,j}^{\frac{1}{p}}}{[\gamma_n(\Omega_n)]^{\frac{1}{p}+r}} \left[ A_{0}\left( f \right) +k_{p}^{\frac{r+1}{p}}A_{r+1}\left( f \right) \right]\\ &\leqslant \max \left\{ \frac{C_{p,r,j}^{\frac{1}{p}}}{[\gamma_n(\Omega_n)]^{\frac{1}{p}+r}} ,\frac{C_{p,r,j}^{\frac{1}{p}}}{[\gamma_n(\Omega_n)]^{\frac{1}{p}+r}} k_{p}^{\frac{r+1}{p}}+k_{p}^{\frac{r+1-j}{p}} \right\} \left( A_{0}\left( f \right) +A_{r+1}\left( f \right) \right) ,\end{aligned}$$
	which yields \eqref{1f1f13f3f1}. This completes the proof of Theorem~\ref{zyao}.
\end{proof}
\section{Sobolev norm equivalence on open convex subsets of \texorpdfstring{$\ell^{2}$}{l2}}
Fix a sequence $\{a_i\}_{i=1}^\infty\subset(0,\infty)$ satisfying $\sum_{i=1}^\infty a_i^2<\infty$, and fix a nonempty open convex set $\Omega\subset\ell^2$.  For $\mathbf x=(x_i)_{i\geqslant1}\in\ell^2$ and $n\in\nn$, write $\mathbf x_n=(x_1,\ldots,x_n)$ for the first $n$ coordinates and $\mathbf x^n=(x_{n+1},x_{n+2},\ldots)$ for the tail.  Define
$$\Omega_{n} \triangleq \left\{ \textbf{x}_{n} =\left( x_{i} \right) \in \mathbb{R}^{n} :\left( a_{1}x_{1},a_{2}x_{2},\cdots ,a_{n}x_{n},\textbf{0}^{n} \right) \in \Omega \right\} .$$
Because $\Omega$ is open and nonempty and the finite-coordinate truncations converge in $\ell^2$, $\Omega_n$ is nonempty for all sufficiently large $n$; convexity and openness follow directly from the definition. 
For later use, set
$$
E_n\triangleq\left\{(x_1,\ldots,x_n)\in\mathbb{R}^n:
(x_1,\ldots,x_n,\mathbf 0^n)\in\Omega\right\}.
$$

For any nonempty set $I\subset \mathbb{N}$, write
\begin{eqnarray*}
	\ell^2(I)\triangleq \left\{\textbf{x}=(x_i)_{i\in I}\in \mathbb{R}^{I}:\sum_{i\in I}|x_i|^2<\infty\right\}.
\end{eqnarray*}
As in Section~2.2 of \cite[pp.~523--525]{YZ}, the product law with one-dimensional marginals $\mathcal N_{a_i}$ is concentrated on $\ell^2$ and defines a Borel probability measure $P$ there. Following \cite[(8), p.~523]{YZ}, for each $k\in\mathbb{N}$, we set $\mathcal{N}^k\triangleq\prod\limits_{i=1}^{k}\mathcal{N}_{a_i}$, where
$$
\mathcal{N}_a(B)\triangleq \frac{1}{\sqrt{2\pi a^2}}\int_Be^{-\frac{x^2}{2a^2}}\mathrm{d}x,\quad\,\forall\, B\in\mathscr{B}(\mathbb{R}).
$$
Since $\ell^2$ can be identified with $\mathbb{R}^k\times \ell^{2}(\mathbb{N}\setminus\{1,2,\ldots,k\})$, we have the decomposition  $P=\mathcal{N}^k\times P^{\widehat{1,2,\ldots,k}}$.  Here $P^{\widehat{1,\ldots,k}}$ denotes the product measure obtained by omitting the $1,2,\ldots,k$-th components; i.e., it is the restriction of the product measure $\prod\limits_{j\in\mathbb{N}\setminus\{1,\ldots,k\}}\mathcal{N}_{a_j}$ to the space
$$
\left(\ell^{2}(\mathbb{N}\setminus\{1,\ldots,k\}),\mathscr{B}\big(\ell^{2}(\mathbb{N}\setminus\{1,\ldots,k\})\big)\right).
$$
\begin{lemma}\label{jihejixian}
	We have
	\begin{eqnarray}
		\lim_{n\rightarrow \infty} \gamma_{n} \left( \Omega_{n} \right) =P\left( \Omega \right) >0.\label{13f13f13g}
	\end{eqnarray}
\end{lemma}
\begin{proof}
	Since every nondegenerate Gaussian measure on a separable Hilbert space is full
	(see, e.g., \cite[Proposition~1.25, p.~21]{DP}), we have $P\left( \Omega \right) >0$.
For $\mathbf{x}=(x_i)_{i\geqslant1}\in\ell^2$, set $\pi_n\mathbf{x}=(x_1,\ldots,x_n,0,0,\ldots)$. Since $\pi_n\mathbf{x}\to\mathbf{x}$ in $\ell^2$, openness of $\Omega$ implies that the following indicator functions converge pointwise to $1$ on $\Omega$ and to $0$ on $\ell^2\setminus\overline{\Omega}$:
	$$\lim_{n\rightarrow \infty} \chi_{E_n \times \ell^{2} \left( \mathbb{N} \setminus \left\{ 1,2,\cdots ,n \right\} \right)} =\begin{cases}1,&\textbf{x} \in \Omega\\ 0,&\textbf{x} \in \ell^{2} \setminus \overline{\Omega}\end{cases}.$$
	Moreover, by the discussion following \cite[Corollary~9, p.~1074]{BBL}, the boundary of a convex set with nonempty interior has Gaussian measure zero. Hence $P\left( \partial \Omega \right) =0$. Therefore, by the dominated convergence theorem, we obtain
	$$\begin{aligned}\lim_{n\rightarrow \infty} \gamma_{n} \left( \Omega_{n} \right)&=\lim_{n\rightarrow \infty} \left( \int_{\left\{ \textbf{x} =\left( x_{i} \right) \in \mathbb{R}^{n} :\left( a_{1}x_{1},a_{2}x_{2},\cdots ,a_{n}x_{n},\textbf{0}^{n} \right) \in \Omega \right\}} 1\mathrm{d} \gamma_{n} \right)\\ &=\lim_{n\rightarrow \infty} \left( \int_{E_n \times \ell^{2} \left( \mathbb{N} \setminus \left\{ 1,2,\cdots ,n \right\} \right)} 1\mathrm{d} P \right)\\ &=\int_{\ell^{2}} \lim_{n\rightarrow \infty} \chi_{E_n \times \ell^{2} \left( \mathbb{N} \setminus \left\{ 1,2,\cdots ,n \right\} \right)} \mathrm{d} P\\ &=\int_{\Omega} 1\mathrm{d}P=P\left( \Omega \right) >0,\end{aligned}$$
	This completes the proof of Lemma~\ref{jihejixian}.
\end{proof}
Recall the function space $C_{0,F^{\infty}}^{\infty}\left( \Omega \right)$ introduced in \cite[p.~7]{WYZ1} (for the real case, see \cite[p.~8]{WYZZ}). By \cite[Lemma~3.1, p.~17]{WYZ1}, we know that $C_{0,F^{\infty}}^{\infty}\left( \Omega \right)$ is dense in $L^{p}\left( \Omega ,P \right)$, and we shall use $C_{0,F^{\infty}}^{\infty}\left( \Omega \right)$ as the test function space. 

\begin{definition}
For $j\in\mathbb{N}$, let  
$$
\delta_j\phi \triangleq \partial_j\phi-\frac{x_j}{a_j^2}\phi,
\qquad \phi\in C_{0,F^\infty}^{\infty}(\Omega).
$$  
Let $k\in\mathbb{N}$, $i_1,\ldots,i_k\in\mathbb{N}$, and $f\in L^1(\Omega,P)$. We say that $f$ has a weak partial derivative of order $k$ in the directions $i_1,\ldots,i_k$ if there exists  
$$
f_{i_1\cdots i_k}\in L^1(\Omega,P)
$$  
such that  
$$
\int_{\Omega}
f\,
\delta_{i_1}\delta_{i_2}\cdots\delta_{i_k}\phi
\,\mathrm dP
=
(-1)^k
\int_{\Omega}
f_{i_1\cdots i_k}\phi
\,\mathrm dP,
\qquad
\forall\phi\in C_{0,F^\infty}^{\infty}(\Omega).
$$  
Here  
$$
\delta_{i_1}\delta_{i_2}\cdots\delta_{i_k}\phi
\triangleq
\delta_{i_1}\bigl(
\delta_{i_2}(\cdots(\delta_{i_k}\phi)\cdots)
\bigr).
$$  
In this case, $f_{i_1\cdots i_k}$ is uniquely determined $P$-almost everywhere, and we write  
$$
\partial_{i_1i_2\cdots i_k}f
\triangleq
f_{i_1\cdots i_k}.
$$

In particular, when $k=1$, this definition reduces to  
$$
\int_{\Omega}
f\,\delta_j\phi\,\mathrm dP
=
-\int_{\Omega}
\partial_j f\,\phi\,\mathrm dP,
\qquad
\forall\phi\in C_{0,F^\infty}^{\infty}(\Omega).
$$

We say that $f$ has weak partial derivatives up to order $m$ if, for every $k=1,\ldots,m$ and every $i_1,\ldots,i_k\in\mathbb{N}$, the weak partial derivative $\partial_{i_1\cdots i_k}f$ exists in the above sense.
\end{definition}

\begin{remark}
The preceding definition is equivalent to the recursive definition  
$$
\int_{\Omega}
\partial_{i_1\cdots i_{k-1}}f\,
\delta_{i_k}\phi
\,\mathrm dP
=
-\int_{\Omega}
\partial_{i_1\cdots i_k}f\,
\phi
\,\mathrm dP,
\qquad
\forall\phi\in C_{0,F^\infty}^{\infty}(\Omega),
$$  
whenever the weak derivatives of orders up to $k$ exist.

Moreover, the operators $\delta_i$ commute on $C_{0,F^\infty}^{\infty}(\Omega)$. Hence, by uniqueness of weak derivatives,  
$$
\partial_{i_1\cdots i_k}f
=
\partial_{i_{\sigma(1)}\cdots i_{\sigma(k)}}f
\qquad P\text{-a.e.}
$$  
for every permutation $\sigma$ of $\{1,\ldots,k\}$.  
Thus the $k$-th weak derivative is symmetric with respect to its indices.
\end{remark}

%By \cite[Theorem 3.2, p.14]{WZ}, we know that every non-empty open subset of $\ell^{2}$ admits a $C_{F^{\infty}}^{\infty}$ exhaustion function. Then, similarly to the proof of (2) and (3) in \cite[Proposition 3.2, pp.24-25]{WYZ1}, we have:
%\begin{lemma}\label{nini1f3ff3}
%	Let $f,\partial_{j}f\in L^{1}\left( \Omega ,loc \right),j\in \nn$. Suppose that for any $\phi \in C_{S}^{1}\left( \Omega \right)$ with $\mathrm{supp} \phi \s \Omega$ and such that $\phi$ and all its partial derivatives are bounded, the following holds:
%	$$\int_{\Omega} f\delta_{j} \phi \mathrm{d} P=-\int_{\Omega} \partial_{j} f\cdot \phi \mathrm{d} P.$$
%\end{lemma}


%The proofs of Lemma \ref{nini1f3ff3} and Lemma \ref{1f3ff32} can be found in \cite{WYZZ}.

\begin{definition}
Let $m\in\nn$. We define
\[
\begin{aligned}
W^{m,p}(\Omega,P)\triangleq\Bigg\{f\in L^p(\Omega,P):\;&
\int_{\Omega}
\left(
\sum_{i_1,\ldots,i_k=1}^{\infty}
\left|a_{i_1}\cdots a_{i_k}\partial_{i_1\cdots i_k}f\right|^2
\right)^{p/2}\mathrm dP<\infty,\\[-1mm]
&k=1,\ldots,m\Bigg\}.
\end{aligned}
\]
For $f\in W^{m,p}(\Omega,P)$, we define
\[
\begin{aligned}
\|f\|_{W^{m,p}(\Omega,P)}
\triangleq {}&
\left(\int_{\Omega}|f|^p\,\mathrm dP\right)^{1/p}\\
&+\sum_{k=1}^{m}
\left[
\int_{\Omega}
\left(
\sum_{i_1,\ldots,i_k=1}^{\infty}
\left|a_{i_1}\cdots a_{i_k}\partial_{i_1\cdots i_k}f\right|^2
\right)^{p/2}\mathrm dP
\right]^{1/p}.
\end{aligned}
\]
\end{definition}
\begin{remark}
		With the weak derivatives defined above, $W^{m,p}(\Omega,P)$ is a Banach space with respect to the norm $\|f\|_{W^{m,p}(\Omega,P)}$.
\end{remark}
Using the same approximation scheme as in the proof of the main result of \cite{WYZZ}, with the weak derivatives defined above, one obtains the following density statement.
\begin{theorem}\label{choumixingdingli}
The restrictions to $\Omega$ of functions in $\mathscr C_P^\infty$ are dense in $W^{m,p}(\Omega,P)$.
	\end{theorem}
\begin{theorem}\label{1ff1f3}
Let $m\geqslant2$.  There exists a constant $K>0$, depending only on $\Omega$, $p$, and $m$, such that, for every $f\in W^{m,p}\left( \Omega ,P\right)$,
\begin{eqnarray}
\begin{aligned}
\|f\|_{W^{m,p}(\Omega,P)}
&\leqslant K\Bigg[
\left(\int_{\Omega}|f|^p\,\mathrm dP\right)^{1/p}\\
&\qquad\quad+
\left(
\int_{\Omega}
\left(
\sum_{i_1,\ldots,i_m=1}^{\infty}
\left|a_{i_1}\cdots a_{i_m}\partial_{i_1\cdots i_m}f\right|^2
\right)^{p/2}\mathrm dP
\right)^{1/p}
\Bigg].
\end{aligned}
\label{13f113vgv3v13v}
\end{eqnarray}
\end{theorem}
\begin{proof}
By Theorem \ref{choumixingdingli}, it suffices to prove, for $f\in C_{P}^{\infty}\left( \mathbb{R}^{N} \right) ,N\in \mathbb{N}$, the inequality
\begin{eqnarray}\label{1ff313g}
\begin{aligned}
&\left(\int_{\Omega}|f|^p\,\mathrm dP\right)^{1/p}
+\sum_{k=1}^{m}
\left[
\int_{\Omega}
\left(
\sum_{i_1,\ldots,i_k=1}^{N}
\left|a_{i_1}\cdots a_{i_k}\partial_{i_1\cdots i_k}f\right|^2
\right)^{p/2}\mathrm dP
\right]^{1/p}
\\
&\qquad\leqslant
K\Bigg\{
\left(\int_{\Omega}|f|^p\,\mathrm dP\right)^{1/p}
+
\left[
\int_{\Omega}
\left(
\sum_{i_1,\ldots,i_m=1}^{N}
\left|a_{i_1}\cdots a_{i_m}\partial_{i_1\cdots i_m}f\right|^2
\right)^{p/2}\mathrm dP
\right]^{1/p}
\Bigg\}.
\end{aligned}
\end{eqnarray}
Consider
$$\begin{aligned}K&\triangleq 1+\sum_{j=1}^{m-1} \max \left\{ \frac{C_{p,m-1,j}^{\frac{1}{p}}}{[P(\Omega)]^{\frac{1}{p}+m-1}} ,\frac{C_{p,m-1,j}^{\frac{1}{p}}}{[P(\Omega)]^{\frac{1}{p}+m-1}} k_{p}^{\frac{m}{p}}+k_{p}^{\frac{m-j}{p}} \right\} ,\\ K_{n}&\triangleq 1+\sum_{j=1}^{m-1} \max \left\{ \frac{C_{p,m-1,j}^{\frac{1}{p}}}{[\gamma_n(\Omega_n)]^{\frac{1}{p}+m-1}} ,\frac{C_{p,m-1,j}^{\frac{1}{p}}}{[\gamma_n(\Omega_n)]^{\frac{1}{p}+m-1}} k_{p}^{\frac{m}{p}}+k_{p}^{\frac{m-j}{p}} \right\} ,n=1,2,\cdots\end{aligned} .$$
By Lemma \ref{jihejixian}, we obtain
\begin{eqnarray}
	\lim_{n\rightarrow \infty} K_{n}=K.\label{1f3ggg}
\end{eqnarray}
For sufficiently large $n>N$ with $\gamma_n(\Omega_n)>0$, regard $f$ as a function on $\rr^n$ independent of the last $n-N$ variables and set $F_n(x_1,\ldots,x_n)\triangleq f(a_1x_1,\ldots,a_Nx_N)$.  Apply inequality \eqref{1f1f13f3f1} of Theorem~\ref{zyao}. Then we have
\begin{eqnarray}
	\sum_{j=0}^{m} A_{j}\left( F_n \right) \leqslant K_{n}\left( A_{0}\left( F_n \right) +A_{m}\left( F_n \right) \right).\label{1f31f1f13}
\end{eqnarray}
By an argument similar to the proof of Lemma \ref{jihejixian}, the dominated convergence theorem yields
\[
\begin{aligned}
\lim_{n\to\infty}A_0(F_n)
&=\lim_{n\to\infty}
\left(\int_{\Omega_n}|F_n|^p\,\mathrm d\gamma_n\right)^{1/p}\\
&=\lim_{n\to\infty}
\left(\int_{E_n}|f|^p\,\mathrm d\mathcal N^n\right)^{1/p}\\
&=\left(\int_{\Omega}|f|^p\,\mathrm dP\right)^{1/p}.
\end{aligned}
\]
and, for each $j=1,2,\cdots,m$, we have
\[
\begin{aligned}
\lim_{n\to\infty}A_j(F_n)
&=\lim_{n\to\infty}
\left[
\int_{\Omega_n}
\left(
\sum_{i_1,\ldots,i_j=1}^{N}
|\partial_{i_1\cdots i_j}F_n|^2
\right)^{p/2}\mathrm d\gamma_n
\right]^{1/p}\\
&=\lim_{n\to\infty}
\left[
\int_{E_n}
\left(
\sum_{i_1,\ldots,i_j=1}^{N}
a_{i_1}^2\cdots a_{i_j}^2
|\partial_{i_1\cdots i_j}f|^2
\right)^{p/2}\mathrm d\mathcal N^n
\right]^{1/p}\\
&=
\left[
\int_{\Omega}
\left(
\sum_{i_1,\ldots,i_j=1}^{N}
a_{i_1}^2\cdots a_{i_j}^2
|\partial_{i_1\cdots i_j}f|^2
\right)^{p/2}\mathrm dP
\right]^{1/p}.
\end{aligned}
\]
Combining \eqref{1f3ggg} with the preceding limits and letting $n\to\infty$ in \eqref{1f31f1f13} gives \eqref{1ff313g}.  This completes the proof of Theorem~\ref{1ff1f3}.
\end{proof}
For $1\leqslant p<2$, the Hilbertian aggregation of coordinate derivatives in the norm above is essential.  If one instead uses the coordinatewise $\ell^p$-type Sobolev norm from \cite{WYZZ}, an analogue of \eqref{13f113vgv3v13v} fails, as the following proposition shows.
\begin{proposition}\label{1f3g}
For $p\in \left[ 1,2 \right)$ and any natural number $m\geqslant 2$, there does not exist a constant $K > 0$ such that, for every $f\in \mathscr{C}^{\infty}_P$,
	\begin{eqnarray}\label{13f311gg}
		\begin{aligned}
			&\int_{\ell^{2}} \left| f \right|^{p} \mathrm{d} P+\sum_{k=1}^{m} \sum_{i_{1},i_{2},\cdots ,i_{k}=1}^{\infty} a_{i_{1}}^{p}a_{i_{2}}^{p}\cdots a_{i_{k}}^{p}\int_{\ell^{2}} \left| \partial_{i_{1}i_{2}\cdots i_{k}} f \right|^{p} \mathrm{d} P\\ &\  \leqslant K\left[ \int_{\ell^{2}} \left| f \right|^{p} \mathrm{d} P+\sum_{i_{1},i_{2},\cdots ,i_{m}=1}^{\infty} a_{i_{1}}^{p}a_{i_{2}}^{p}\cdots a_{i_{m}}^{p}\int_{\ell^{2}} \left| \partial_{i_{1}i_{2}\cdots i_{m}} f \right|^{p} \mathrm{d} P \right]\end{aligned}
	\end{eqnarray}
\end{proposition}
\begin{proof}
	Suppose that there exists a constant $K > 0$ such that \eqref{13f311gg} holds for every $f\in \mathscr{C}^{\infty}_P$. Then for every $n\in\nn$, $f=\sum_{i=1}^{n} \frac{x_{i}}{a_{i}}$ should satisfy inequality \eqref{13f311gg}.
	
	A direct computation shows that inequality \eqref{13f311gg} becomes
	$$\frac{2^{\frac{p}{2}}\Gamma \left( \frac{p+1}{2} \right)}{\sqrt{\pi}} n^{\frac{p}{2}}+n\leqslant K\frac{2^{\frac{p}{2}}\Gamma \left( \frac{p+1}{2} \right)}{\sqrt{\pi}} n^{\frac{p}{2}},n=1,2,\cdots,$$
	which is clearly impossible. This contradiction proves Proposition~\ref{1f3g}.
\end{proof}
\section{Malliavin--Sobolev norm equivalence}
\label{sec:malliavin-L1-open-problem}

We now apply the preceding dimension-free estimates to Malliavin spaces.  The scalar case follows by representing a cylindrical random variable on a finite-dimensional Gaussian space.  The Hilbert-space-valued case is then obtained by finite-rank reduction and Gaussian randomization, in the spirit of \cite{AddonaMuratoriRossi}.

The constants obtained in Section~3 are independent of the dimension of the Gaussian representation.  This allows the estimates to pass to an arbitrary isonormal Gaussian process and gives the equivalence between the full Sobolev norm and the graph norm of the highest Malliavin derivative for every $p\in[1,+\infty)$ and every integer $k\geqslant2$.

We use a different notation from \cite{AddonaMuratoriRossi}.  In \cite{AddonaMuratoriRossi}, the letter $q$ denotes the integrability exponent and the letter $p$ denotes the order of the Malliavin derivative.  In the present paper we use $p$ for the integrability exponent and $k$ for the order of differentiation.  Thus, the case $p=1$ and $k\geqslant3$ below is exactly the case $q=1$ and derivative order at least three left open in \cite{AddonaMuratoriRossi}.



Let $(\Xi,\mathscr{F},\mathbb{P})$ be a probability space, let $(\mathfrak{H},\langle\cdot,\cdot\rangle_{\mathfrak{H}})$ be a real separable Hilbert space, and let
$$
W=\{W(h):h\in\mathfrak{H}\}
$$
be an isonormal Gaussian process over $\mathfrak{H}$, namely
$$
\mathbb{E}[W(h)W(g)]=\langle h,g\rangle_{\mathfrak{H}},\qquad h,g\in\mathfrak{H}.
$$
Let $\mathcal{S}$ denote the class of smooth cylindrical random variables
$$
F=f\bigl(W(h_{1}),\ldots,W(h_{m})\bigr),
$$
where $m\in\mathbb{N}$, $h_{1},\ldots,h_{m}\in\mathfrak{H}$ and $f\in C_{P}^{\infty}(\mathbb{R}^{m})$.  For $j\in\mathbb{N}$ we denote by $D^{j}F$ the $j$-th Malliavin derivative, regarded as an $\mathfrak{H}^{\otimes j}$-valued random variable.

For $p\in[1,+\infty)$ and $k\geqslant2$, define on $\mathcal{S}$ the full Sobolev norm
$$
\|F\|_{\mathcal{D}(k,p)}
\triangleq
\|F\|_{L^{p}(\Xi)}
+\sum_{j=1}^{k}\|D^{j}F\|_{L^{p}(\Xi;\mathfrak{H}^{\otimes j})},
$$
and the graph norm of $D^{k}$
$$
\|F\|_{\mathcal{G}(k,p)}
\triangleq
\|F\|_{L^{p}(\Xi)}
+\|D^{k}F\|_{L^{p}(\Xi;\mathfrak{H}^{\otimes k})}.
$$
Clearly,
$$
\|F\|_{\mathcal{G}(k,p)}\leqslant \|F\|_{\mathcal{D}(k,p)}.
$$
The nontrivial issue is the reverse inequality with a constant independent of the dimension of the cylindrical representation of $F$.

We first isolate the finite-dimensional consequence of Theorem~\ref{zyao}.

\begin{lemma}[Dimension-free scalar Gaussian interpolation]
	\label{lem:scalar-dimension-free-gaussian}
	Let $p\in[1,+\infty)$ and let $k\geqslant2$.  For every $j\in\{1,\ldots,k-1\}$ there exists a constant $L_{p,k,j}>0$, depending only on $p,k,j$, such that for every $n\in\mathbb{N}$ and every $f\in C_{P}^{\infty}(\mathbb{R}^{n})$,
	$$
	\left(\int_{\mathbb{R}^{n}}\|D^{j}f\|_{j}^{p}\,\mathrm{d}\gamma_{n}\right)^{\frac1p}
	\leqslant
	L_{p,k,j}
	\left[
	\left(\int_{\mathbb{R}^{n}}|f|^{p}\,\mathrm{d}\gamma_{n}\right)^{\frac1p}
	+
	\left(\int_{\mathbb{R}^{n}}\|D^{k}f\|_{k}^{p}\,\mathrm{d}\gamma_{n}\right)^{\frac1p}
	\right].
	$$
	In particular, if
	$$
	L_{p,k}\triangleq\sum_{j=1}^{k-1}L_{p,k,j},
	$$
	then
	\[
\begin{aligned}
&\left(\int_{\mathbb{R}^{n}}|f|^{p}\,\mathrm{d}\gamma_{n}\right)^{1/p}
+
\sum_{j=1}^{k}
\left(\int_{\mathbb{R}^{n}}\|D^{j}f\|_{j}^{p}\,\mathrm{d}\gamma_{n}\right)^{1/p}
\\
&\qquad\leqslant
(1+L_{p,k})
\left[
\left(\int_{\mathbb{R}^{n}}|f|^{p}\,\mathrm{d}\gamma_{n}\right)^{1/p}
+
\left(\int_{\mathbb{R}^{n}}\|D^{k}f\|_{k}^{p}\,\mathrm{d}\gamma_{n}\right)^{1/p}
\right].
\end{aligned}
\]
	The constants are independent of $n$.
\end{lemma}

\begin{proof}
	Apply Theorem~\ref{zyao} with $\Omega_{n}=\mathbb{R}^{n}$ and $r=k-1$.  Since $\gamma_{n}(\mathbb{R}^{n})=1$, the proof of Theorem~\ref{zyao} gives, for $j=1,\ldots,k-1$,
	$$
	A_{j}(f)
	\leqslant
	L_{p,k,j}\bigl(A_{0}(f)+A_{k}(f)\bigr),
	$$
	where one may take
	$$
	L_{p,k,j}
	\triangleq
	\max\left\{
	C_{p,k-1,j}^{\frac1p},
	C_{p,k-1,j}^{\frac1p}k_{p}^{\frac{k}{p}}
	+k_{p}^{\frac{k-j}{p}}
	\right\}.
	$$
	By the definition of $A_{j}$, this is exactly the first displayed estimate.  Summing over $j=1,\ldots,k-1$ and adding the zeroth- and $k$-th-order terms proves the second one.  The constants $C_{p,k-1,j}$ and $k_p$ are independent of $n$; hence the same is true of $L_{p,k,j}$ and $L_{p,k}$.  This completes the proof of Lemma~\ref{lem:scalar-dimension-free-gaussian}.
\end{proof}

We next pass from the finite-dimensional lemma to the scalar Malliavin setting.

\begin{theorem}[Scalar Malliavin--Sobolev norm equivalence]
	\label{thm:scalar-malliavin-equivalence-all-p-k}
	Let $p\in[1,+\infty)$ and $k\geqslant2$.  Then, for every real separable Hilbert space $\mathfrak{H}$, every isonormal Gaussian process $W$ over $\mathfrak{H}$ and every $F\in\mathcal{S}$,
	$$
	\|F\|_{\mathcal{G}(k,p)}
	\leqslant
	\|F\|_{\mathcal{D}(k,p)}
	\leqslant
	(1+L_{p,k})\|F\|_{\mathcal{G}(k,p)}.
	$$
	In particular, the equivalence constant depends only on $p$ and $k$ and is independent of $\dim\mathfrak{H}$.
\end{theorem}

\begin{proof}
	Only the second inequality needs to be proved.  Let
	$$
	F=f\bigl(W(h_{1}),\ldots,W(h_{m})\bigr)\in\mathcal{S}.
	$$
	Set
	$$
	E\triangleq\operatorname{span}\{h_{1},\ldots,h_{m}\}
	$$
	and let $n=\dim E$.  Choose an orthonormal basis $e_{1},\ldots,e_{n}$ of $E$.  Since each $h_{\ell}$ is a linear combination of $e_{1},\ldots,e_{n}$, there exists $g\in C_{P}^{\infty}(\mathbb{R}^{n})$ such that
	$$
	F=g\bigl(W(e_{1}),\ldots,W(e_{n})\bigr).
	$$
	Because $e_{1},\ldots,e_{n}$ are orthonormal, the random vector
	$$
	\bigl(W(e_{1}),\ldots,W(e_{n})\bigr)
	$$
	has law $\gamma_{n}$.  Moreover, for every $j\in\{1,\ldots,k\}$,
	$$
	D^{j}F
	=
	\sum_{i_{1},\ldots,i_{j}=1}^{n}
	\partial_{i_{1}\cdots i_{j}}g\bigl(W(e_{1}),\ldots,W(e_{n})\bigr)
	\,e_{i_{1}}\otimes\cdots\otimes e_{i_{j}}.
	$$
	Therefore,
	$$
	\|D^{j}F\|_{L^{p}(\Xi;\mathfrak{H}^{\otimes j})}
	=
	\left(
	\int_{\mathbb{R}^{n}}\|D^{j}g\|_{j}^{p}\,\mathrm{d}\gamma_{n}
	\right)^{\frac1p},
	$$
	and similarly
	$$
	\|F\|_{L^{p}(\Xi)}
	=
	\left(
	\int_{\mathbb{R}^{n}}|g|^{p}\,\mathrm{d}\gamma_{n}
	\right)^{\frac1p}.
	$$
	Applying Lemma~\ref{lem:scalar-dimension-free-gaussian} to $g$ gives
	$$
	\sum_{j=1}^{k-1}
	\|D^{j}F\|_{L^{p}(\Xi;\mathfrak{H}^{\otimes j})}
	\leqslant
	L_{p,k}
	\left(
	\|F\|_{L^{p}(\Xi)}
	+
	\|D^{k}F\|_{L^{p}(\Xi;\mathfrak{H}^{\otimes k})}
	\right).
	$$
	Adding the zeroth- and $k$-th-order terms proves Theorem~\ref{thm:scalar-malliavin-equivalence-all-p-k}.
\end{proof}




We now extend Theorem~\ref{thm:scalar-malliavin-equivalence-all-p-k} to random variables with values in an arbitrary real separable Hilbert space.  This is the form naturally considered in \cite{AddonaMuratoriRossi}.

Let $(V,\langle\cdot,\cdot\rangle_{V})$ be a real separable Hilbert space.  Denote by $\mathcal{S}_{V}$ the class of smooth finite-rank $V$-valued random variables
$$
F=\sum_{\ell=1}^{M}F_{\ell}v_{\ell},
\qquad
F_{\ell}\in\mathcal{S},\quad v_{\ell}\in V.
$$
For $F\in\mathcal{S}_{V}$ we define
$$
D^{j}F
\triangleq
\sum_{\ell=1}^{M}D^{j}F_{\ell}\otimes v_{\ell}
\in
L^{p}\bigl(\Xi;\mathfrak{H}^{\otimes j}\otimes V\bigr).
$$
The corresponding full and graph norms are
$$
\|F\|_{\mathcal{D}(k,p)(V)}
\triangleq
\|F\|_{L^{p}(\Xi;V)}
+
\sum_{j=1}^{k}
\|D^{j}F\|_{L^{p}(\Xi;\mathfrak{H}^{\otimes j}\otimes V)},
$$
and
$$
\|F\|_{\mathcal{G}(k,p)(V)}
\triangleq
\|F\|_{L^{p}(\Xi;V)}
+
\|D^{k}F\|_{L^{p}(\Xi;\mathfrak{H}^{\otimes k}\otimes V)}.
$$

We shall use the following dimension-free consequence of the Gaussian
Kahane--Khintchine inequality; see, for instance,
\cite[Theorem~6.2.6, p.~23]{HNVW}.

\begin{lemma}
	\label{lem:gaussian-randomization-hilbert}
There exist constants
	$0<\alpha_p\leqslant\beta_p<+\infty$, depending only on $p$, such that
	for every real Hilbert space $U$, every $M\in\mathbb{N}$, and every
	$u_1,\ldots,u_M\in U$,
	\[
	\alpha_p
	\left(\sum_{\ell=1}^{M}\|u_\ell\|_U^2\right)^{\frac12}
	\leqslant
	\left(
	\int_{\mathbb{R}^M}
	\left\|\sum_{\ell=1}^{M}y_\ell u_\ell\right\|_U^p
	\,\mathrm{d}\gamma_M(\mathbf{y})
	\right)^{\frac1p}
	\leqslant
	\beta_p
	\left(\sum_{\ell=1}^{M}\|u_\ell\|_U^2\right)^{\frac12}.
	\]
	For $p=2$, one may take $\alpha_2=\beta_2=1$.
\end{lemma}

\begin{proof}
	By the Gaussian Kahane--Khintchine inequality
\cite[Theorem~6.2.6, p.~23]{HNVW},
	for every $p,q\in[1,+\infty)$ there exists a constant
	$C_{p,q}>0$, depending only on $p$ and $q$, such that
\begin{eqnarray}\label{wahha}
		\left(
	\int_{\mathbb{R}^M}
	\left\|\sum_{\ell=1}^{M}y_\ell u_\ell\right\|_U^p
	\,\mathrm{d}\gamma_M(\mathbf{y})
	\right)^{\frac1p}
	\leqslant
	C_{p,q}
	\left(
	\int_{\mathbb{R}^M}
	\left\|\sum_{\ell=1}^{M}y_\ell u_\ell\right\|_U^q
	\,\mathrm{d}\gamma_M(\mathbf{y})
	\right)^{\frac1q}.
\end{eqnarray}
	The constant $C_{p,q}$ is independent of $M$, $U$, and the vectors
	$u_1,\ldots,u_M$.
	
	Taking $q=2$ and using that the coordinate functions
	$y_1,\ldots,y_M$ are independent standard Gaussian random variables,
	we obtain
\begin{eqnarray}\label{1f1ff}
		\begin{aligned}
		\int_{\mathbb{R}^M}
		\left\|\sum_{\ell=1}^{M}y_\ell u_\ell\right\|_U^2
		\,\mathrm{d}\gamma_M(\mathbf{y})
		&=
		\sum_{\ell,k=1}^{M}
		\left(
		\int_{\mathbb{R}^M}y_\ell y_k
		\,\mathrm{d}\gamma_M(\mathbf{y})
		\right)
		\langle u_\ell,u_k\rangle_U
		\\
		&=
		\sum_{\ell=1}^{M}\|u_\ell\|_U^2.
	\end{aligned}
\end{eqnarray}
	Applying  \eqref{wahha} with the pairs
	$(p,2)$ and $(2,p)$ therefore yields the claimed estimate, for instance
	with
	\[
	\alpha_p=C_{2,p}^{-1},
	\qquad
	\beta_p=C_{p,2}.
	\]
	When $p=2$, the identity \eqref{1f1ff} gives the assertion with
	$\alpha_2=\beta_2=1$.  This completes the proof of Lemma~\ref{lem:gaussian-randomization-hilbert}.
\end{proof}

We now turn to the Hilbert-space-valued estimate.

\begin{theorem}[Hilbert-valued Malliavin--Sobolev norm equivalence]
	\label{thm:vector-malliavin-equivalence-all-p-k}
	Let $p\in[1,+\infty)$ and $k\geqslant2$.  Let $\mathfrak{H}$ and $V$ be arbitrary real separable Hilbert spaces and let $W$ be an isonormal Gaussian process over $\mathfrak{H}$.  Then, for every $F\in\mathcal{S}_{V}$,
	$$
	\|F\|_{\mathcal{G}(k,p)(V)}
	\leqslant
	\|F\|_{\mathcal{D}(k,p)(V)}
	\leqslant
	K_{p,k}\|F\|_{\mathcal{G}(k,p)(V)},
	$$
	where
	$$
	K_{p,k}
	\triangleq
	1+\frac{\beta_{p}}{\alpha_{p}}L_{p,k}.
	$$
	In particular, $K_{p,k}$ depends only on $p$ and $k$ and is independent of $\dim\mathfrak{H}$ and $\dim V$.
\end{theorem}

\begin{proof}
	The first inequality is immediate.  We prove the reverse one.
	
	Let $F\in\mathcal{S}_{V}$.  Since $F$ has finite-dimensional range, after replacing the original spanning family in $V$ by an orthonormal basis of its span, we may write
	$$
	F=\sum_{\ell=1}^{M}F_{\ell}v_{\ell},
	$$
	where $v_{1},\ldots,v_{M}$ are orthonormal in $V$ and $F_{1},\ldots,F_{M}\in\mathcal{S}$.
	
	For $\textbf{y}=(y_{1},\ldots,y_{M})\in\mathbb{R}^{M}$ define the scalar-valued smooth random variable
	$$
	F_{\textbf{y}}
	\triangleq
	\sum_{\ell=1}^{M}y_{\ell}F_{\ell}.
	$$
	For every $j\geqslant1$,
	$$
	D^{j}F_{\textbf{y}}
	=
	\sum_{\ell=1}^{M}y_{\ell}D^{j}F_{\ell}.
	$$
	Fix $j\in\{1,\ldots,k-1\}$.  By the scalar estimate in the proof of Theorem~\ref{thm:scalar-malliavin-equivalence-all-p-k}, for every $\textbf{y}\in\mathbb{R}^{M}$,
	$$
	\|D^{j}F_{\textbf{y}}\|_{L^{p}(\Xi;\mathfrak{H}^{\otimes j})}
	\leqslant
	L_{p,k,j}
	\left[
	\|F_{\textbf{y}}\|_{L^{p}(\Xi)}
	+
	\|D^{k}F_{\textbf{y}}\|_{L^{p}(\Xi;\mathfrak{H}^{\otimes k})}
	\right].
	$$
	Taking the $L^{p}(\mathbb{R}^{M},\gamma_{M})$ norm in the parameter $\textbf{y}$ and using Minkowski's inequality, we get
	$$
	\begin{aligned}
		&\left(
		\int_{\mathbb{R}^{M}}
		\|D^{j}F_{\textbf{y}}\|_{L^{p}(\Xi;\mathfrak{H}^{\otimes j})}^{p}
		\,\mathrm{d}\gamma_{M}(\textbf{y})
		\right)^{\frac1p}
		\\
		&\quad\leqslant
		L_{p,k,j}
		\left(
		\int_{\mathbb{R}^{M}}
		\|F_{\textbf{y}}\|_{L^{p}(\Xi)}^{p}
		\,\mathrm{d}\gamma_{M}(\textbf{y})
		\right)^{\frac1p}
		\\
		&\qquad\quad+
		L_{p,k,j}
		\left(
		\int_{\mathbb{R}^{M}}
		\|D^{k}F_{\textbf{y}}\|_{L^{p}(\Xi;\mathfrak{H}^{\otimes k})}^{p}
		\,\mathrm{d}\gamma_{M}(\textbf{y})
		\right)^{\frac1p}.
	\end{aligned}
	$$
	By Fubini's theorem and Lemma~\ref{lem:gaussian-randomization-hilbert}, applied pointwise in $\omega\in\Xi$, the left-hand side is bounded from below by
	$$
	\alpha_{p}
	\|D^{j}F\|_{L^{p}(\Xi;\mathfrak{H}^{\otimes j}\otimes V)}.
	$$
	Indeed, since $v_{1},\ldots,v_{M}$ are orthonormal,
	$$
	\sum_{\ell=1}^{M}
	\|D^{j}F_{\ell}(\omega)\|_{\mathfrak{H}^{\otimes j}}^{2}
	=
	\|D^{j}F(\omega)\|_{\mathfrak{H}^{\otimes j}\otimes V}^{2}.
	$$
	Similarly, the two terms on the right-hand side are bounded from above by
	$$
	\beta_{p}\|F\|_{L^{p}(\Xi;V)}
	$$
	and
	$$
	\beta_{p}
	\|D^{k}F\|_{L^{p}(\Xi;\mathfrak{H}^{\otimes k}\otimes V)},
	$$
	respectively.  Hence
	$$
	\|D^{j}F\|_{L^{p}(\Xi;\mathfrak{H}^{\otimes j}\otimes V)}
	\leqslant
	\frac{\beta_{p}}{\alpha_{p}}L_{p,k,j}
	\left[
	\|F\|_{L^{p}(\Xi;V)}
	+
	\|D^{k}F\|_{L^{p}(\Xi;\mathfrak{H}^{\otimes k}\otimes V)}
	\right].
	$$
	Summing this estimate over $j=1,\ldots,k-1$ gives
	$$
	\sum_{j=1}^{k-1}
	\|D^{j}F\|_{L^{p}(\Xi;\mathfrak{H}^{\otimes j}\otimes V)}
	\leqslant
	\frac{\beta_{p}}{\alpha_{p}}L_{p,k}
	\|F\|_{\mathcal{G}(k,p)(V)}.
	$$
	Adding the zeroth- and $k$-th-order terms yields
	$$
	\|F\|_{\mathcal{D}(k,p)(V)}
	\leqslant
	\left(1+\frac{\beta_{p}}{\alpha_{p}}L_{p,k}\right)
	\|F\|_{\mathcal{G}(k,p)(V)},
	$$
	which proves the theorem.
\end{proof}

The norm equivalence also identifies the corresponding completions.  Let $\mathbb{D}^{k,p}(V)$ denote the completion of $\mathcal{S}_{V}$ with respect to $\|\cdot\|_{\mathcal{D}(k,p)(V)}$, and let $\mathbb{D}_{*}^{k,p}(V)$ denote the completion of $\mathcal{S}_{V}$ with respect to $\|\cdot\|_{\mathcal{G}(k,p)(V)}$.

\begin{corollary}
	\label{cor:malliavin-domains-coincide-all-p-k}
	Let $p\in[1,+\infty)$ and $k\geqslant2$.  For arbitrary real separable Hilbert spaces $\mathfrak{H}$ and $V$,
	$$
	\mathbb{D}^{k,p}(V)=\mathbb{D}_{*}^{k,p}(V)
	$$
	with equivalent norms.  More precisely, the identity map on $\mathcal{S}_{V}$ extends uniquely to an isomorphism between the two completions, and
	$$
	\|F\|_{\mathcal{G}(k,p)(V)}
	\leqslant
	\|F\|_{\mathcal{D}(k,p)(V)}
	\leqslant
	K_{p,k}\|F\|_{\mathcal{G}(k,p)(V)}.
	$$
\end{corollary}

\begin{proof}
	Theorem~\ref{thm:vector-malliavin-equivalence-all-p-k} shows that the two norms are equivalent on the common dense subspace $\mathcal{S}_{V}$.  Equivalent norms have canonically isomorphic completions, which proves Corollary~\ref{cor:malliavin-domains-coincide-all-p-k}.
\end{proof}

\begin{remark}[Unified range of exponents and orders]
	The argument above gives one proof covering simultaneously all
	$$
	p\in[1,+\infty),\qquad k\geqslant2.
	$$
	For $p>1$ it recovers the previously known norm equivalence in infinite-dimensional Malliavin spaces.  For $p=1$ and $k=2$ it recovers the endpoint result obtained in \cite{AddonaMuratoriRossi}.  The genuinely new regime is $p=1$ and $k\geqslant3$, which is exactly the open case identified there.
\end{remark}


\section*{Acknowledgements}

\begin{thebibliography}{99}
\bibitem{AdamsFournier}
R.~A.~Adams and J.~J.~F.~Fournier,
\emph{Sobolev Spaces}, 2nd ed.,
Pure and Applied Mathematics, Vol.~140,
Elsevier/Academic Press, Amsterdam, 2003.

\bibitem{AddonaWeighted2025}
D.~Addona,
\emph{Equivalence of Sobolev norms with respect to weighted Gaussian
	measures},
{\sl Potential Anal.}, \textbf{62} (2025), 761--816.

\bibitem{AddonaMuratoriRossi}
D.~Addona, M.~Muratori and M.~Rossi,
\emph{On the equivalence of Sobolev norms in Malliavin spaces},
{\sl J. Funct. Anal.}, \textbf{283} (2022), Paper No.~109600.

\bibitem{AngiuliFerrariPallara2019}
L.~Angiuli, S.~Ferrari and D.~Pallara,
\emph{Gradient estimates for perturbed Ornstein--Uhlenbeck semigroups
	on infinite-dimensional convex domains},
{\sl J. Evol. Equ.}, \textbf{19} (2019), 677--715.

\bibitem{AS17}
G.~Aubrun and S.~J.~Szarek,
\emph{Alice and Bob Meet Banach: The Interface of Asymptotic Geometric
	Analysis and Quantum Information Theory},
Math. Surveys Monogr., Vol.~223,
American Mathematical Society, Providence, RI, 2017.

\bibitem{BBL}
J.~M.~Borwein, J.~V.~Burke and A.~S.~Lewis,
\emph{Differentiability of cone-monotone functions on separable Banach
	space},
{\sl Proc. Amer. Math. Soc.}, \textbf{132} (2004), 1067--1076.

\bibitem{BogachevGaussian}
V.~I.~Bogachev,
\emph{Gaussian Measures},
Mathematical Surveys and Monographs, Vol.~62,
American Mathematical Society, Providence, RI, 1998.

\bibitem{BogachevSurvey}
V.~I.~Bogachev,
\emph{Sobolev classes on infinite-dimensional spaces},
in: \emph{Geometric Measure Theory and Real Analysis},
Publications of the Scuola Normale Superiore, Vol.~17,
Edizioni della Normale, Pisa, 2014, 1--56.

\bibitem{BogachevPilipenkoShaposhnikov}
V.~I.~Bogachev, A.~Yu.~Pilipenko and A.~V.~Shaposhnikov,
\emph{Sobolev functions on infinite-dimensional domains},
{\sl J. Math. Anal. Appl.}, \textbf{419} (2014), 1023--1044.

\bibitem{CappaFerrari2016}
G.~Cappa and S.~Ferrari,
\emph{Maximal Sobolev regularity for solutions of elliptic equations
	in infinite dimensional Banach spaces endowed with a weighted
	Gaussian measure},
{\sl J. Differential Equations}, \textbf{261} (2016), 7099--7131.

\bibitem{CappaFerrari2018}
G.~Cappa and S.~Ferrari,
\emph{Maximal Sobolev regularity for solutions of elliptic equations
	in Banach spaces endowed with a weighted Gaussian measure:
	the convex subset case},
{\sl J. Math. Anal. Appl.}, \textbf{458} (2018), 300--331.

\bibitem{CW}
A.~Carbery and J.~Wright,
\emph{Distributional and $L^q$ norm inequalities for polynomials over
	convex bodies in $\mathbb R^n$},
{\sl Math. Res. Lett.}, \textbf{8} (2001), 233--248.

\bibitem{Daletskii1967}
Yu.~L.~Daletskii,
\emph{Infinite-dimensional elliptic operators and corresponding
	parabolic equations},
{\sl Russian Math. Surveys}, \textbf{22} (1967), 1--53.

\bibitem{DP}
G.~Da Prato,
\emph{An Introduction to Infinite-Dimensional Analysis},
Springer-Verlag, Berlin, 2006.

\bibitem{DaPratoLunardi2014}
G.~Da Prato and A.~Lunardi,
\emph{Sobolev regularity for a class of second order elliptic PDE's in
	infinite dimension},
{\sl Ann. Probab.}, \textbf{42} (2014), 2113--2160.

\bibitem{DaPratoLunardi2015}
G.~Da Prato and A.~Lunardi,
\emph{Maximal Sobolev regularity in Neumann problems for gradient
	systems in infinite dimensional domains},
{\sl Ann. Inst. Henri Poincar\'e Probab. Stat.},
\textbf{51} (2015), 1102--1123.

\bibitem{EI}
A.~Eskenazis and P.~Ivanisvili,
\emph{Dimension independent Bernstein--Markov inequalities in Gauss
	space},
{\sl J. Approx. Theory}, \textbf{253} (2020), Paper No.~105377.

\bibitem{Ferrari2019}
S.~Ferrari,
\emph{Sobolev spaces with respect to a weighted Gaussian measure in
	infinite dimensions},
{\sl Infin. Dimens. Anal. Quantum Probab. Relat. Top.},
\textbf{22} (2019), Paper No.~1950026.

\bibitem{FeyelPradelle1989}
D.~Feyel and A.~de La Pradelle,
\emph{Espaces de Sobolev gaussiens},
{\sl Ann. Inst. Fourier (Grenoble)}, \textbf{39} (1989), 875--908.

\bibitem{Frolov1973}
N.~N.~Frolov,
\emph{On a coercive inequality for an elliptic operator in infinitely
	many independent variables},
{\sl Math. USSR-Sb.}, \textbf{19} (1973), 395--406.

\bibitem{GozlanJuillet}
N.~Gozlan and N.~Juillet,
\emph{On a mixture of Brenier and Strassen theorems},
{\sl Proc. Lond. Math. Soc.}, \textbf{120} (2020), 434--463.

\bibitem{GrossAWS}
L.~Gross,
\emph{Abstract Wiener spaces},
in: \emph{Proceedings of the Fifth Berkeley Symposium on Mathematical
	Statistics and Probability, Vol. II, Part 1},
University of California Press, Berkeley, 1967, 31--42.

\bibitem{GrossPotential}
L.~Gross,
\emph{Potential theory on Hilbert space},
{\sl J. Funct. Anal.}, \textbf{1} (1967), 123--181.

\bibitem{Gross}
L.~Gross,
\emph{Logarithmic Sobolev inequalities},
{\sl Amer. J. Math.}, \textbf{97} (1975), 1061--1083.

\bibitem{Gundy1989}
R.~F.~Gundy,
\emph{Some Topics in Probability and Analysis},
CBMS Regional Conference Series in Mathematics, Vol.~70,
American Mathematical Society, Providence, RI, 1989.

\bibitem{Hino2003}
M.~Hino,
\emph{On Dirichlet spaces over convex sets in infinite dimensions},
in: \emph{Finite and Infinite Dimensional Analysis in Honor of
	Leonard Gross},
Contemp. Math., Vol.~317,
American Mathematical Society, Providence, RI, 2003, 143--156.

\bibitem{Hino2011}
M.~Hino,
\emph{Dirichlet spaces on $H$-convex sets in Wiener space},
{\sl Bull. Sci. Math.}, \textbf{135} (2011), 667--683.

\bibitem{HNVW}
T.~Hyt\"onen, J.~van Neerven, M.~Veraar and L.~Weis,
\emph{Analysis in Banach Spaces. Vol. II:
	Probabilistic Methods and Operator Theory},
Ergebnisse der Mathematik und ihrer Grenzgebiete, Vol.~67,
Springer, Cham, 2017.

\bibitem{Kree1977}
M.~Kr\'ee,
\emph{Propri\'et\'e de trace en dimension infinie d'espaces du type
	Sobolev},
{\sl Bull. Soc. Math. France}, \textbf{105} (1977), 141--163.

\bibitem{LedouxOleszkiewicz}
M.~Ledoux and K.~Oleszkiewicz,
\emph{On measure concentration of vector-valued maps},
{\sl Bull. Polish Acad. Sci. Math.}, \textbf{55} (2007), 261--278.

\bibitem{LQ17}
D.~Li and H.~Queff\'elec,
\emph{Introduction to Banach Spaces: Analysis and Probability},
Vol.~1, Cambridge University Press, Cambridge, 2017.

\bibitem{Malliavin1978}
P.~Malliavin,
\emph{Stochastic calculus of variations and hypoelliptic operators},
in: \emph{Proceedings of the International Symposium on Stochastic
	Differential Equations, Kyoto, 1976},
Wiley, New York, 1978, 195--263.

\bibitem{Meyer1982}
P.-A.~Meyer,
\emph{Note sur les processus d'Ornstein--Uhlenbeck},
in: \emph{S\'eminaire de Probabilit\'es XVI},
Lecture Notes in Math., Vol.~920,
Springer, Berlin, 1982, 95--132.

\bibitem{Meyer1983}
P.-A.~Meyer,
\emph{Quelques r\'esultats analytiques sur le semi-groupe
	d'Ornstein--Uhlenbeck en dimension infinie},
in: \emph{Theory and Application of Random Fields},
Lecture Notes in Control and Information Sciences, Vol.~49,
Springer, Berlin, 1983, 201--214.

\bibitem{Meyer1984}
P.-A.~Meyer,
\emph{Transformations de Riesz pour les lois gaussiennes},
in: \emph{S\'eminaire de Probabilit\'es XVIII},
Lecture Notes in Math., Vol.~1059,
Springer, Berlin, 1984, 179--193.

\bibitem{Nualart2006}
D.~Nualart,
\emph{The Malliavin Calculus and Related Topics},
2nd ed., Probability and its Applications,
Springer-Verlag, Berlin, 2006.

\bibitem{Pisier1988}
G.~Pisier,
\emph{Riesz transforms: a simpler analytic proof of P.~A.~Meyer's
	inequality},
in: \emph{S\'eminaire de Probabilit\'es XXII},
Lecture Notes in Math., Vol.~1321,
Springer, Berlin, 1988, 485--501.

\bibitem{PronkVeraar}
M.~Pronk and M.~Veraar,
\emph{Tools for Malliavin calculus in UMD Banach spaces},
{\sl Potential Anal.}, \textbf{40} (2014), 307--344.

\bibitem{Shigekawa1980}
I.~Shigekawa,
\emph{Derivatives of Wiener functionals and absolute continuity of
	induced measures},
{\sl J. Math. Kyoto Univ.}, \textbf{20} (1980), 263--289.

\bibitem{ShigekawaBook}
I.~Shigekawa,
\emph{Stochastic Analysis},
Translations of Mathematical Monographs, Vol.~224,
American Mathematical Society, Providence, RI, 2004.

\bibitem{Stroock1981}
D.~W.~Stroock,
\emph{The Malliavin calculus, a functional analytic approach},
{\sl J. Funct. Anal.}, \textbf{44} (1981), 212--257.

\bibitem{Sugita1985}
H.~Sugita,
\emph{Sobolev spaces of Wiener functionals and Malliavin's calculus},
{\sl J. Math. Kyoto Univ.}, \textbf{25} (1985), 31--48.

\bibitem{SugitaChar1985}
H.~Sugita,
\emph{On a characterization of the Sobolev spaces over an abstract
	Wiener space},
{\sl J. Math. Kyoto Univ.}, \textbf{25} (1985), 717--725.

\bibitem{Ustunel1995}
A.~S.~\"Ust\"unel,
\emph{An Introduction to Analysis on Wiener Space},
Lecture Notes in Mathematics, Vol.~1610,
Springer-Verlag, Berlin, 1995.

\bibitem{Watanabe1984}
S.~Watanabe,
\emph{Lectures on Stochastic Differential Equations and Malliavin
	Calculus},
Tata Institute of Fundamental Research Lectures on Mathematics and
Physics, Vol.~73,
Springer-Verlag, Berlin, 1984.

\bibitem{YZ}
J.~Yu and X.~Zhang,
\emph{$L^2$ estimates and existence theorems for the
	$\overline{\partial}$ operators in infinite dimensions, I},
{\sl J. Math. Pures Appl.}, \textbf{163} (2022), 518--548.

\bibitem{WYZ1}
Z.~Wang, J.~Yu and X.~Zhang,
\emph{$L^2$ estimates and existence theorems for the
	$\overline{\partial}$ operators in infinite dimensions, II},
{\sl J. Math. Pures Appl.}, \textbf{205} (2026),
Paper No.~103811.

\bibitem{WYZZ}
Z.~Wang, J.~Yu, X.~Zhang and S.~Zhao,
\emph{``$H=W$'' in infinite dimensions},
arXiv:2602.04136, version 2, 2026.
\end{thebibliography}

\end{document}
